\chapter{Spaces of Constant Curvature}
\label{chap:dc-spaces-of-constant-curvature}

Multiplying a Riemannian metric by $c>0$ divides its sectional curvature by
$c$. Thus constant-curvature metrics may be normalized to curvature $1$, $0$,
or $-1$. The corresponding simply connected models are $S^n$, $\mathbb R^n$,
and $H^n$.

\section{Theorem of Cartan on the determination of the metric by means of the curvature}

Let $M$ and $\tilde{M}$ be two Riemannian manifolds of dimension $n$ and let
$p\in M$ and $\tilde{p}\in\tilde{M}$. Choose a linear isometry
$i:T_p(M)\to T_{\tilde{p}}(\tilde{M})$. Let $V\subset M$ be a normal neighborhood
of $p$ such that $\exp_{\tilde{p}}$ is defined at $i\circ\exp_p^{-1}(V)$. Define a
mapping $f:V\to\tilde{M}$ by

$$
f(q)=\exp_{\tilde{p}}\circ\,i\circ\exp_p^{-1}(q),\qquad q\in V.
$$

For $q\in V$, let $\gamma:[0,t]\to M$
with $\gamma(0)=p$, $\gamma(t)=q$. Denote by $P_t$ the parallel transport along
$\gamma$ from $\gamma(0)$ to $\gamma(t)$. Define $\phi_t:T_q(M)\to T_{f(q)}(\tilde{M})$
by

$$
\phi_t(v)=\tilde{P}_t\circ i\circ P_t^{-1}(v),\qquad v\in T_q(M),
$$

where $\tilde{P}_t$ is the parallel transport along the normalized geodesic
$\tilde{\gamma}:[0,t]\to\tilde{M}$ given by $\tilde{\gamma}(0)=\tilde{p}$,
$\tilde{\gamma}'(0)=i(\gamma'(0))$. Denote by $R$ and $\tilde{R}$ the curvatures of
$M$ and $\tilde{M}$, respectively.

\begin{lemma}[Jacobi transfer in parallel frames]
  \label{lem:dc-ch8-2-1-jacobi-transfer}
  \lean{Riemannian.Jacobi.jacobiFrameTransfer}
  \uses{def:dc-ch5-2-1}
  \leanok
  Let $\gamma$ and $\tilde\gamma$ be geodesics with parallel orthonormal frames
  $(e_i)$ and $(\tilde e_i)$, where $e_n=\gamma'$ and
  $\tilde e_n=\tilde\gamma'$. Suppose that $J=\sum_i y_i e_i$ satisfies the Jacobi
  equation $\tfrac{D^2 J}{dt^2}+R(u',J)u'=0$ along $\gamma$, and that at the parameter
  $t$ the two frames have matching curvature coefficients,

  $$
  \langle R(e_n,e_i)e_n,e_j\rangle=\langle\tilde R(\tilde e_n,\tilde e_i)\tilde e_n,\tilde e_j\rangle,
  \qquad i,j=1,\dots,n.
  $$

  Then $\tilde J=\sum_i y_i\tilde e_i$ is a Jacobi field along $\tilde\gamma$.
\end{lemma}
\begin{proof}
  \leanok
\sketch
  Parallelism reduces the Jacobi equation to
  $y_j''+\sum_i a_{ij}y_i=0$, where
  $a_{ij}=\langle R(e_n,e_i)e_n,e_j\rangle$. The matching hypothesis gives the
  identical system for $(\tilde e_i)$, and reassembling its solution yields
  $D^2\tilde J/dt^2+\tilde R(\tilde\gamma',\tilde J)\tilde\gamma'=0$.
\end{proof}

\begin{lemma}[Jacobi transfer in constant curvature]
  \label{lem:dc-ch8-2-1-jacobi-transfer-const}
  \lean{Riemannian.Jacobi.jacobiFrameTransfer_isConstantCurvature}
  \uses{lem:dc-ch8-2-1-jacobi-transfer}
  \leanok
  Suppose $M$ and $\tilde M$ both have constant sectional curvature $K_0$ (the same value),
  and let $e_1,\dots,e_n$ resp.\ $\tilde e_1,\dots,\tilde e_n$ be parallel orthonormal frames
  along geodesics $\gamma$ resp.\ $\tilde\gamma$, with the velocities $u'=e_{n_0}$,
  $\tilde u'=\tilde e_{n_0}$ among the frame vectors. Then the curvature-matching hypothesis
  of \cref{lem:dc-ch8-2-1-jacobi-transfer},

  $$
  \langle R(e_i,u')u',e_j\rangle=\langle\tilde R(\tilde e_i,\tilde u')\tilde u',\tilde e_j\rangle,
  $$

  holds automatically: the constant-curvature identity and orthonormality make both sides
  $K_0(\delta_{ij}-\delta_{i n_0}\delta_{n_0 j})$. Hence a Jacobi field $J=\sum_i y_i e_i$ along
  $\gamma$ carries over, with the \emph{same} scalar coefficients, to a Jacobi field
  $\tilde J=\sum_i y_i\tilde e_i$ along $\tilde\gamma$.
\end{lemma}
\begin{proof}
  \leanok
\sketch
  The constant-curvature identity gives
  $\langle R(a,u')u',b\rangle=K_0(\langle u',u'\rangle\langle a,b\rangle-\langle a,u'\rangle\langle u',b\rangle)$
  Substituting
  $a=e_i$, $b=e_j$, $u'=e_{n_0}$ and using orthonormality
  ($\langle e_i,e_j\rangle=\delta_{ij}$, $\langle e_i,e_{n_0}\rangle=\delta_{i n_0}$,
  $\langle e_{n_0},e_{n_0}\rangle=1$) collapses each entry to
  $K_0(\delta_{ij}-\delta_{i n_0}\delta_{n_0 j})$. The two frame-coefficient
  matrices therefore coincide, discharging the hypothesis of
  \cref{lem:dc-ch8-2-1-jacobi-transfer}, whose conclusion is the claimed carry-over.
\end{proof}

\begin{lemma}[velocity-seeded parallel orthonormal frame]
  \label{lem:dc-ch8-2-1-velocity-frame}
  \lean{Riemannian.Jacobi.exists_velocitySeededParallelOrthoFrame}
  \uses{def:dc-ch5-2-1}
  \leanok
  Let $\gamma:[a,b]\to M$ be a geodesic whose initial velocity is a unit vector. Then there is a parallel
  orthonormal frame $e_1(t),\dots,e_n(t)$ along $\gamma$ (parallel: $De_i/dt=0$; orthonormal:
  $\langle e_i(t),e_j(t)\rangle=\delta_{ij}$ for all $t$) whose distinguished member is the
  velocity, $e_{n_0}(t)=\gamma'(t)$ for all $t\in[a,b]$.
\end{lemma}
\begin{proof}
  \leanok
\sketch
  Extend $\gamma'(a)$ to an orthonormal basis and parallel-transport that basis.
  Parallel transport preserves the metric, while uniqueness identifies the transported
  copy of $\gamma'(a)$ with the parallel field $\gamma'$.
\end{proof}

\begin{lemma}[the parallel orthonormal frame along a geodesic, intrinsically]
  \label{lem:dc-ch8-2-1-parallel-ortho-frame}
  \lean{Riemannian.Jacobi.exists_parallelOrthoFrameAlongOn}
  \uses{def:dc-ch5-2-1}
  \leanok
  Let $\gamma:[a,b]\to M$ be a geodesic and let $\{e^0_i\}$ be an orthonormal family in
  $T_{\gamma(a)}M$. Then there is a family $\{e_i(t)\}$ of vector fields along $\gamma$ with
  $e_i(a)=e^0_i$, each $e_i$ parallel along $\gamma$, and
  $\langle e_i(t),e_j(t)\rangle_{\gamma(t)}=\delta_{ij}$ for every $t\in[a,b]$.
\end{lemma}
\begin{proof}
  \leanok
\sketch
  Parallel-transport each $e^0_i$ along $\gamma$; this exists with any prescribed initial
  value. The pairing of two parallel fields is constant, so the Gram matrix remains
  equal to $(\delta_{ij})$.
\end{proof}

\begin{lemma}[Parallel frame transported by a linear isometry]
  \label{lem:dc-ch8-2-1-transported-frame}
  \lean{Riemannian.Jacobi.exists_transportedParallelOrthoFrame, Riemannian.Jacobi.exists_transportedParallelOrthoFrame_pair}
  \uses{lem:dc-ch8-2-1-parallel-ortho-frame}
  Let $i:T_pM\to T_{\tilde p}\tilde M$ be a linear isometry, let $\{e^0_j\}$ be an orthonormal
  family at $p=\gamma(a)$, and let $\tilde\gamma:[a,b]\to\tilde M$ be a geodesic with
  $\tilde\gamma(a)=\tilde p$. Then there is a family $\{\tilde e_j(t)\}$ along $\tilde\gamma$
  with $\tilde e_j(a)=i(e^0_j)$, each $\tilde e_j$ parallel along $\tilde\gamma$, and
  $\langle\tilde e_j(t),\tilde e_k(t)\rangle_{\tilde\gamma(t)}=\delta_{jk}$ for every
  $t\in[a,b]$. If $e_j$ is the frame from the corresponding initial data along
  $\gamma$, then $\tilde e_j(t)=\phi_t(e_j(t))$ for
  $\phi_t=\tilde P_t\circ i\circ P_t^{-1}$.
\end{lemma}
\begin{proof}
  \leanok
\sketch
  The family $(i(e_j^0))$ is orthonormal. Apply
  \cref{lem:dc-ch8-2-1-parallel-ortho-frame} and use uniqueness of parallel
  transport together with \cref{lem:dc-ch8-2-1-phi}.
\end{proof}

\begin{lemma}[uniqueness of parallel transport, and the transport map $P_t$]
  \label{lem:dc-ch8-2-1-parallel-transport-map}
  \lean{Riemannian.Jacobi.IsParallelFieldAlongOn.eqOn_of_initial, Riemannian.Jacobi.parallelTransportAlong, Riemannian.Jacobi.metricInner_parallelTransportAlong, Riemannian.Jacobi.parallelTransportAlongEquiv}
  \uses{def:dc-ch5-2-1}
  \leanok
  Let $\gamma:[a,b]\to M$ be a geodesic. Two fields parallel along $\gamma$ that agree at $a$
  agree on all of $[a,b]$. Consequently the assignment carrying $w_0\in T_{\gamma(a)}M$ to the
  value at $t$ of the parallel field along $\gamma$ seeded by $w_0$ is a well-defined map

  $$
  P_t:T_{\gamma(a)}M\longrightarrow T_{\gamma(t)}M,\qquad t\in[a,b],
  $$

  which is linear, satisfies $P_a=\mathrm{id}$, and is an isometry:
  $\langle P_tu,P_tw\rangle_{\gamma(t)}=\langle u,w\rangle_{\gamma(a)}$. In particular $P_t$ is
  injective, hence a linear isomorphism onto $T_{\gamma(t)}M$.
\end{lemma}
\begin{proof}
  \leanok
  \uses{def:dc-ch5-2-1}
\sketch
  Let $v,w$ be parallel along $\gamma$ with $v(a)=w(a)$, and set $d=v-w$. The parallel
  transport equation $\frac{Dd}{dt}=0$ is linear in $d$, so $d$ is parallel along $\gamma$,
  and $d(a)=0$. The pairing of two parallel fields along $\gamma$ is constant, so

  $$
  \langle d(t),d(t)\rangle_{\gamma(t)}=\langle d(a),d(a)\rangle_{\gamma(a)}=0
  \qquad(t\in[a,b]),
  $$

  and since $g_{\gamma(t)}$ is positive definite, $d(t)=0$; that is, $v(t)=w(t)$.

  Uniqueness makes $P_t$ well defined, and existence of a parallel field with prescribed value
  at $a$ makes it total. For linearity, let $v,w$ be the parallel fields seeded by $w_0,w_1$.
  Then $v+w$ is parallel and $(v+w)(a)=w_0+w_1$, so by uniqueness $v+w$ is the parallel field
  seeded by $w_0+w_1$; evaluating at $t$ gives $P_t(w_0+w_1)=P_tw_0+P_tw_1$, and homogeneity is
  the same argument applied to $rv$. $P_a=\mathrm{id}$ is the seeding condition. The isometry
  statement is once more the constancy of the pairing of two parallel fields, applied to the
  fields seeded by $u$ and $w$. Finally, if $P_tw_0=0$ then
  $\langle w_0,w_0\rangle_{\gamma(a)}=\langle P_tw_0,P_tw_0\rangle_{\gamma(t)}=0$, so $w_0=0$;
  an injective endomorphism of a finite-dimensional space is an isomorphism.
\end{proof}

\begin{lemma}[Parallel-transport conjugation]
  \label{lem:dc-ch8-2-1-phi}
  \lean{Riemannian.Jacobi.parallelTransportConjugate, Riemannian.Jacobi.parallelTransportConjugate_left, Riemannian.Jacobi.metricInner_parallelTransportConjugate, Riemannian.Jacobi.eq_parallelTransportConjugate_of_isParallelFieldAlongOn}
  \uses{lem:dc-ch8-2-1-parallel-transport-map}
  \leanok
  Let $\gamma:[a,b]\to M$ and $\tilde\gamma:[a,b]\to\tilde M$ be geodesics and let
  $i:T_{\gamma(a)}M\to T_{\tilde\gamma(a)}\tilde M$ be linear. For $t\in[a,b]$ put

  $$
  \phi_t=\tilde P_t\circ i\circ P_t^{-1}:
  T_{\gamma(t)}M\longrightarrow T_{\tilde\gamma(t)}\tilde M,
  $$

  with $P_t$, $\tilde P_t$ the parallel transports of
  \cref{lem:dc-ch8-2-1-parallel-transport-map}. Then $\phi_a=i$; if $i$ is a linear isometry
  then so is every $\phi_t$; and if $v$ is parallel along $\gamma$ and $\tilde v$ is parallel
  along $\tilde\gamma$ with $\tilde v(a)=i(v(a))$, then

  $$
  \tilde v(t)=\phi_t(v(t))\qquad\text{for every }t\in[a,b].
  $$

\end{lemma}
\begin{proof}
  \leanok
  \uses{lem:dc-ch8-2-1-parallel-transport-map}
\sketch
  $\phi_a=i$ because $P_a$ and $\tilde P_a$ are the identity. If $i$ is a linear isometry then
  $\phi_t$ is a composite of three isometries, $P_t^{-1}$ and $\tilde P_t$ being isometries by
  \cref{lem:dc-ch8-2-1-parallel-transport-map}; hence $\phi_t$ is one.

  For the last claim, $v$ is parallel with value $v(a)$ at $a$, so by the uniqueness clause of
  \cref{lem:dc-ch8-2-1-parallel-transport-map} it is the parallel field seeded by $v(a)$, i.e.
  $P_t(v(a))=v(t)$, whence $P_t^{-1}(v(t))=v(a)$. Therefore

  $$
  \phi_t(v(t))=\tilde P_t\big(i(v(a))\big)=\tilde P_t\big(\tilde v(a)\big)=\tilde v(t),
  $$

  the last equality because $\tilde v$ is parallel along $\tilde\gamma$ with value $\tilde v(a)$
  at $a$, hence is the parallel field seeded by $\tilde v(a)$.
\end{proof}

\begin{lemma}[the chart frame coefficient is the intrinsic curvature form]
  \label{lem:dc-ch8-2-1-curvature-bridge}
  \lean{Riemannian.Jacobi.chartMetricInner_chartCurvatureEndo_eq_curvatureFormAt, Riemannian.Jacobi.chartFrameRealize_tangentCoordChange, Riemannian.Jacobi.chartMetricInner_chartCurvatureEndo_chartVectorRep_eq_curvatureFormAt}
  \leanok
  For $v,a,b\in T_qM$, the frame coefficient $\langle\mathcal R(a,v)v,b\rangle$ computed in a
  chart --- the form in which the Jacobi equation is actually run --- equals the intrinsic
  curvature form $R(v,a,v,b)$ evaluated on $v,a,v,b$ themselves. No curvature hypothesis is
  assumed.
\end{lemma}
\begin{proof}
  \leanok
\sketch
  Pairing the chart Jacobi operator against a chart vector gives
  $\langle\mathcal R(a,v)v,b\rangle$. The manifold--chart curvature bridge identifies
  $R(\hat a,\hat v,\hat v,\hat b)$ --- the intrinsic form on the chart-frame realizations
  $\hat\cdot$ of the chart vectors --- with the \emph{negative} of that quantity, and
  antisymmetry of $R$ in its first pair rewrites $R(\hat a,\hat v,\hat v,\hat b)$ as
  $-R(\hat v,\hat a,\hat v,\hat b)$; the two signs cancel. Finally, realizing the chart
  reading of a tangent vector in the chart frame at its own foot returns the vector
  (the readback inverse to the chart reading), so for intrinsic $v,a,b$ no realizations
  remain.
\end{proof}

\begin{lemma}[curvature intertwining in frame coefficients]
  \label{lem:dc-ch8-2-1-hmatch-transfer}
  \lean{Riemannian.Jacobi.chartMetricInner_chartCurvatureEndo_transfer_of_curvatureFormAt}
  \uses{lem:dc-ch8-2-1-curvature-bridge}
  \leanok
  Let $\varphi:T_qM\to T_{\tilde q}\tilde M$ be any map under which the curvature forms
  correspond,

  $$
  \langle R(x,y)z,w\rangle=\big\langle\tilde R(\varphi x,\varphi y)\varphi z,\varphi w\big\rangle
  \qquad\text{for all }x,y,z,w\in T_qM ,
  $$

  which is the hypothesis of \cref{thm:dc-ch8-2-1} with $\varphi=\phi_t$. Then the chart frame
  coefficients of $M$ and $\tilde M$ agree: for all $v,a,b\in T_qM$,

  $$
  \langle\mathcal R(a,v)v,b\rangle=\big\langle\tilde{\mathcal R}(\varphi a,\varphi v)\varphi v,
  \varphi b\big\rangle .
  $$

  Taking $v=\gamma'$ and $a,b$ frame vectors, this is the intrinsic content of the matching
  hypothesis of \cref{lem:dc-ch8-2-1-jacobi-transfer} and of
  \cref{lem:dc-ch8-2-1-jacobi-norm-transfer-general} --- the single point at which the
  curvature hypothesis of \cref{thm:dc-ch8-2-1} is consumed.

  It is not yet that matching hypothesis itself. Those lemmas read
  their frames in a fixed chart along a coordinate curve, so identifying the two demands three
  further things, none of them about curvature: that $\varphi(\gamma'(t))=\tilde\gamma'(t)$
  (which needs $\tilde\gamma'(a)=i(\gamma'(a))$); that its frame vectors are the chart
  readings of the present intrinsic ones under the present $\varphi$; and that its
  parallelism hypothesis, stated in one fixed chart, follows from parallelism along
  $\gamma$ in the sense of \cref{lem:dc-ch8-2-1-parallel-ortho-frame}, which is chart-local.
  That interface is the unbuilt residual of \cref{thm:dc-ch8-2-1}.

  Note $\varphi$ is arbitrary: no linearity, continuity or isometry is used, and no
  parallel-transport theory enters --- $\varphi$ is data, supplied by \cref{thm:dc-ch8-2-1}'s
  own hypothesis as $\phi_t$. The additional structure of $\phi_t$ is what makes the
  transported frame orthonormal (\cref{lem:dc-ch8-2-1-transported-frame}), which is a separate
  requirement.
\end{lemma}
\begin{proof}
  \leanok
\sketch
  Rewrite each side by \cref{lem:dc-ch8-2-1-curvature-bridge} into the intrinsic curvature
  forms $R(v,a,v,b)$ and $\tilde R(\varphi v,\varphi a,\varphi v,\varphi b)$, and apply the
  hypothesis at $(v,a,v,b)$.
\end{proof}

\begin{lemma}[the Jacobi field read in a parallel orthonormal frame]
  \label{lem:dc-ch8-2-1-frame-coeff-bridge}
  \lean{Riemannian.Jacobi.isJacobiPairOn_of_isJacobiFieldOn}
  \uses{def:dc-ch5-2-1, lem:dc-ch8-2-1-parallel-ortho-frame}
  \leanok
  Let $\gamma$ be a geodesic of $M$, read in a fixed chart as the coordinate curve $u$
  staying in the chart domain, and let the pair $(J,DJ)$ satisfy the Jacobi pair system
  along $u$ on a parameter window $[a',b']$,

  $$
  \tfrac{D}{dt}J=DJ,\qquad \tfrac{D}{dt}DJ=-\mathcal R(J,u')u' ,
  $$

  where $\langle\cdot,\cdot\rangle$ and $\mathcal R$ denote the metric and curvature read in
  that chart. Let $e_1(t),\dots,e_n(t)$ ($n=\dim M$) be a parallel orthonormal frame along
  $u$ in the same chart. For a parameter interval $[a,b]\subset(a',b')$ set

  $$
  F_i(t)=\langle J(t),e_i(t)\rangle,\qquad V_i(t)=\langle DJ(t),e_i(t)\rangle .
  $$

  Then on $[a,b]$ the pair $(F,V)$ solves the first-order system

  $$
  F'=V,\qquad V'=-A(t)F,\qquad A(t)_{ij}=a_{ij}(t)=\langle\mathcal R(e_i,u')u',e_j\rangle ,
  $$

  i.e.\ $F_j''+\sum_i a_{ij}F_i=0$.
\end{lemma}
\begin{proof}
  \leanok
\sketch
  Metric compatibility and $De_j/dt=0$ give
  $\frac d{dt}\langle X,e_j\rangle=\langle DX/dt,e_j\rangle$.
  Apply this first to $J$ and then to $DJ$. Expanding $J=\sum_iF_ie_i$ in the
  second identity gives $F'=V$ and $V'=-AF$.
\end{proof}

\begin{lemma}[transfer of the Jacobi norm under matching curvature coefficients]
  \label{lem:dc-ch8-2-1-jacobi-norm-transfer-general}
  \lean{Riemannian.Jacobi.chartMetricInner_transfer_of_jacobiCoef_match}
  \uses{lem:dc-ch8-2-1-frame-coeff-bridge, lem:dc-ch8-2-1-hmatch-transfer, def:dc-ch5-2-1}
  \leanok
  Let $\gamma$ in $M$ and $\tilde\gamma$ in $\tilde M$ ($\dim M=\dim\tilde M=n$) be geodesics
  read in fixed charts as the coordinate curves $u$ and $\tilde u$, each staying in its chart
  domain, let $(J,DJ)$ and $(\tilde J,D\tilde J)$ satisfy the Jacobi pair systems along $u$
  and $\tilde u$ on a window $[a',b']$, and let $e$ and $\tilde e$ be parallel orthonormal
  frames along $u$ and $\tilde u$ in those charts; as in
  \cref{lem:dc-ch8-2-1-frame-coeff-bridge}, $\langle\cdot,\cdot\rangle$ and $\mathcal R$
  denote the metric and curvature read in those charts. Let $[a,b]\subset(a',b')$, and
  suppose that on $[a,b]$ the chart curvature operator $w\mapsto\mathcal R(w,u')u'$ along
  $u$, the frame vectors $e_i$, and the chart metric coefficients along $u$ are all
  continuous, and that the curvature coefficients agree there,

  $$
  \langle\mathcal R(e_i,u')u',e_j\rangle
  =\big\langle\tilde{\mathcal R}(\tilde e_i,\tilde u')\tilde u',\tilde e_j\big\rangle
  \qquad\text{for all }i,j\text{ and all }t\in[a,b],
  $$

  and that the two fields have the same frame data at $a$,
  $\langle J(a),e_i(a)\rangle=\langle\tilde J(a),\tilde e_i(a)\rangle$ and
  $\langle DJ(a),e_i(a)\rangle=\langle D\tilde J(a),\tilde e_i(a)\rangle$ for all $i$.
  Then

  $$
  |\tilde J(t)|^2=|J(t)|^2\qquad\text{for all }t\in[a,b] .
  $$

\end{lemma}
\begin{proof}
  \leanok
\sketch
  By \cref{lem:dc-ch8-2-1-frame-coeff-bridge}, both coefficient pairs solve
  $F'=V$, $V'=-A(t)F$ with the same continuous matrix $A$. Their initial data agree,
  so uniqueness gives $F=\tilde F$.

  Finally, orthonormality of each frame reads the norm off the coefficients: expanding
  $J=\sum_i F_ie_i$ and $\tilde J=\sum_i\tilde F_i\tilde e_i$,

  $$
  |J(t)|^2=\sum_i F_i(t)^2,\qquad |\tilde J(t)|^2=\sum_i\tilde F_i(t)^2 ,
  $$

  and these agree because $F=\tilde F$ on $[a,b]$.
\end{proof}

\begin{lemma}[the transported frame and velocity, read in a fixed chart]
  \label{lem:dc-ch8-2-1-frame-chart-interface}
  \lean{Riemannian.Jacobi.jacobiCoef_match_of_curvatureFormAt, Riemannian.Jacobi.jacobiCoef_match_of_curvatureFormAt_parallelTransportConjugate}
  \uses{lem:dc-ch8-2-1-hmatch-transfer, lem:dc-ch8-2-1-transported-frame, lem:dc-ch8-2-1-parallel-transport-map, lem:dc-ch8-2-1-phi}
  \leanok
  Let $\gamma:[a,b]\to M$ and $\tilde\gamma:[a,b]\to\tilde M$ be geodesics, each of which
  \emph{stays in the source of one fixed chart}, read there as the coordinate curves $u$ and
  $\tilde u$. Let $e_1,\dots,e_n$ be parallel fields along $\gamma$ and $\tilde e_1,\dots,
  \tilde e_n$ parallel fields along $\tilde\gamma$ with matched seeds
  $\tilde e_j(a)=i(e_j(a))$, and suppose $\tilde\gamma'(a)=i(\gamma'(a))$; read all of them
  in those same charts. Then the curvature hypothesis of \cref{thm:dc-ch8-2-1} implies the
  matching hypothesis of \cref{lem:dc-ch8-2-1-jacobi-norm-transfer-general} in the form that
  lemma requires it, namely between the \emph{chart} frame coefficients:

  $$
  \langle\mathcal R(e_i,u')u',e_j\rangle
  =\big\langle\tilde{\mathcal R}(\tilde e_i,\tilde u')\tilde u',\tilde e_j\big\rangle
  \qquad\text{for all }i,j\text{ and all }t\in[a,b] .
  $$

  More generally, the same conclusion holds with $\phi_t$ replaced by any family of
  maps $T_{\gamma(t)}M\to T_{\tilde\gamma(t)}\tilde M$ under which the curvature forms
  correspond and which carries $\gamma'(t)$ to $\tilde\gamma'(t)$ and $e_j(t)$ to
  $\tilde e_j(t)$: no linearity, continuity or isometry of the family is used, and no
  parallel-transport theory enters. The frames need not contain the velocity or be
  orthonormal for this coefficient identity.
\end{lemma}
\begin{proof}
  \leanok
\sketch
  Since the velocity fields are parallel,

  $$
  \phi_t(\gamma'(t))=\tilde P_t\big(i(P_t^{-1}\gamma'(t))\big)=\tilde P_t\big(i(\gamma'(a))\big)
  =\tilde P_t(\tilde\gamma'(a))=\tilde\gamma'(t) ,
  $$

  using $\tilde\gamma'(a)=i(\gamma'(a))$. Also
  $\tilde e_j(t)=\phi_t(e_j(t))$ by \cref{lem:dc-ch8-2-1-transported-frame}.
  Apply \cref{lem:dc-ch8-2-1-hmatch-transfer} with $\varphi=\phi_t$ at
  $(v,a,b)=(\gamma'(t),e_i(t),e_j(t))$ gives the matching of the chart frame coefficients of
  the \emph{chart readings} of $\gamma'(t),e_i(t),e_j(t)$ and of their images under $\phi_t$,
  which gives the required chart coefficient identity.
\end{proof}

\begin{lemma}[the Cartan Jacobi norm transfer in variable curvature, intrinsically]
  \label{lem:dc-ch8-2-1-jacobi-norm-transfer-manifold}
  \lean{Riemannian.Jacobi.metricInner_jacobiField_transfer_of_curvatureFormAt}
  \uses{lem:dc-ch8-2-1-frame-chart-interface, lem:dc-ch8-2-1-jacobi-norm-transfer-general, lem:dc-ch8-2-1-parallel-ortho-frame, def:dc-ch5-2-1}
  \leanok
  Let $\gamma,\tilde\gamma:[a',b']\to M,\tilde M$ be geodesics, each staying in the source of
  one fixed chart; let $e_1,\dots,e_n$ and $\tilde e_1,\dots,\tilde e_n$ be parallel
  \emph{orthonormal} frames along them; and let $\phi_t$ carry $\gamma'(t)$ to
  $\tilde\gamma'(t)$, $e_j(t)$ to $\tilde e_j(t)$, and the curvature form of $M$ at
  $\gamma(t)$ to that of $\tilde M$ at $\tilde\gamma(t)$. Let $J$ be
  a Jacobi field along $\gamma$ and $\tilde J$ one along $\tilde\gamma$, and let
  $[a,b]\subset(a',b')$. If the two have the same frame data at $a$, then

  $$
  |\tilde J(t)|=|J(t)|\qquad\text{for all }t\in[a,b] .
  $$

  No constancy or sign condition on the curvature is required.
\end{lemma}
\begin{proof}
  \leanok
\sketch
  Read the fields and frames in the fixed charts. Their chart readings preserve norms and
  satisfy the hypotheses of \cref{lem:dc-ch8-2-1-jacobi-norm-transfer-general}; coefficient
  matching is \cref{lem:dc-ch8-2-1-frame-chart-interface}. Translate the resulting norm
  equality back to the tangent bundles.
\end{proof}

\begin{lemma}[time reversal of the Jacobi system]
  \label{lem:dc-ch8-2-1-jacobi-reversal}
  \lean{Riemannian.Jacobi.IsJacobiFieldAlongOn.comp_neg}
  \uses{def:dc-ch5-2-1}
  \leanok
  Let $\gamma$ be a geodesic on $[a,b]$ and let $(J,DJ)$ be a Jacobi pair along it. Then
  $t\mapsto\gamma(-t)$ is a geodesic on $[-b,-a]$, and

  $$
  J^-(t)=J(-t),\qquad DJ^-(t)=-DJ(-t)
  $$

  is a Jacobi pair along it.
\end{lemma}
\begin{proof}
  \leanok
\sketch
  Substitute $t\mapsto-t$ into the chart Jacobi system. Bilinearity supplies the signs in
  the connection terms, while the two reversed velocity slots leave the curvature term
  unchanged.
\end{proof}

\begin{lemma}[backward uniqueness of Jacobi fields]
  \label{lem:dc-ch8-2-1-jacobi-backward-uniqueness}
  \lean{Riemannian.Jacobi.IsJacobiFieldAlongOn.eqOn_zero_of_right}
  \uses{lem:dc-ch8-2-1-jacobi-reversal, def:dc-ch5-2-1}
  \leanok
  A Jacobi field along a geodesic on $[a,b]$ with $J(b)=0$ and $DJ(b)=0$ vanishes identically
  on $[a,b]$.
\end{lemma}
\begin{proof}
  \leanok
\sketch
  Reverse time using \cref{lem:dc-ch8-2-1-jacobi-reversal}. The reversed pair has zero
  initial data at $-b$, so forward uniqueness makes it, and hence the original pair, zero.
\end{proof}

\begin{lemma}[a Jacobi field with prescribed data at an interior time]
  \label{lem:dc-ch8-2-1-jacobi-interior-data}
  \lean{Riemannian.Jacobi.exists_isJacobiFieldAlongOn_at}
  \uses{lem:dc-ch8-2-1-jacobi-backward-uniqueness, def:dc-ch5-2-1}
  \leanok
  Let $\gamma$ be a geodesic on $[a,b]$, let $t_0\in[a,b]$, and let $J_0,DJ_0\in
  T_{\gamma(t_0)}M$. Then there is a Jacobi field $(J,DJ)$ on $[a,b]$ with
  $J(t_0)=J_0$ and $DJ(t_0)=DJ_0$.
\end{lemma}
\begin{proof}
  \leanok
\sketch
  Send initial data at $a$ to the resulting data at $t_0$. This linear map is injective by
  \cref{lem:dc-ch8-2-1-jacobi-backward-uniqueness}, and its domain and codomain both have
  dimension $2n$. It is therefore surjective.
\end{proof}

\begin{lemma}[from the norm transfer to the local isometry, in variable curvature]
  \label{lem:dc-ch8-2-1-exp-norm-transfer-general}
  \lean{Riemannian.Jacobi.metricInner_mfderiv_eq_of_semiconjugacy_of_curvatureFormAt}
  \uses{lem:dc-ch8-2-1-jacobi-norm-transfer-manifold, cor:dc-ch5-2-5, lem:dc-ch8-2-1-jacobi-interior-data, lem:dc-ch8-2-1-exp-norm-transfer, lem:dc-ch8-2-1-polarization, lem:dc-ch8-2-1-exp-equiv, lem:dc-ch8-2-1-mfderiv-bridge, lem:dc-ch8-2-2-semiconjugacy, lem:dc-ch8-2-1-no-conjugate}
  \leanok
  Let $M,\tilde M$ be complete, $p\in M$, $\tilde p\in\tilde M$, let $i:T_pM\to
  T_{\tilde p}\tilde M$ be a linear isometry, and let $v\in T_pM$. Write $\gamma=\gamma_v$
  and $\tilde\gamma=\gamma_{i(v)}$ for the two geodesics, and $q=\exp_p(v)$. Suppose:
  \begin{itemize}
    \item \emph{(single chart)} for some window $[a',b']$ with $a'<0<1<b'$, each of $\gamma$,
      $\tilde\gamma$ stays in the source of \emph{one} chart on $[a',b']$;
    \item \emph{(frames)} $e_1,\dots,e_n$ and $\tilde e_1,\dots,\tilde e_n$ are parallel
      \emph{orthonormal} frames along $\gamma$, $\tilde\gamma$ on $[a',b']$;
    \item \emph{(curvature intertwining)} $\phi_t$ carries $\gamma'(t)$ to
      $\tilde\gamma'(t)$, $e_j(t)$ to $\tilde e_j(t)$, and the curvature form of $M$ at
      $\gamma(t)$ to that of $\tilde M$ at $\tilde\gamma(t)$, for $t\in[a',b']$; and
      $\phi_0=i$;
    \item \emph{(no conjugacy)} $1$ is not a conjugate parameter along $\gamma$ or along
      $\tilde\gamma$.
  \end{itemize}
  Let $f$ be \emph{any} map differentiable at $q$ that satisfies the semiconjugacy
  $f\circ\exp_p=\exp_{\tilde p}\circ\,i$ near $v$. Then $f$ preserves the metric at $q$:

  $$
  \langle u,u'\rangle_q=\langle df_q(u),df_q(u')\rangle_{f(q)}\qquad\text{for all }u,u'\in T_qM .
  $$

  This is the variable-curvature counterpart of \cref{lem:dc-ch8-2-2-semiconjugacy}.
\end{lemma}
\begin{proof}
  \leanok
\sketch
  Since $1$ is not conjugate along $\gamma$, $d(\exp_p)_v$ is onto
  (\cref{lem:dc-ch8-2-1-exp-equiv}), so it suffices to prove the claim on vectors
  $d(\exp_p)_v(Z)$. Realize each side as the endpoint of a Jacobi field vanishing at $0$:
  by \cref{cor:dc-ch5-2-5}, $d(\exp_p)_v(Z)=J(1)$ where $J(0)=0$, $DJ(0)=Z$, and likewise
  $d(\exp_{\tilde p})_{i(v)}(i(Z))=\tilde J(1)$ with
  $\tilde J(0)=0$, $D\tilde J(0)=i(Z)$. The fields on the full window are supplied by
  \cref{lem:dc-ch8-2-1-jacobi-interior-data}.

  The two fields have matching frame data at $0$. For the value this is trivial, both being
  $0$. For the covariant derivative, $\tilde e_k(0)=\phi_0(e_k(0))=i(e_k(0))$, so

  $$
  \langle i(Z),\tilde e_k(0)\rangle_{\tilde p}=\langle i(Z),i(e_k(0))\rangle_{\tilde p}
    =\langle Z,e_k(0)\rangle_p
  $$

  because $i$ is a linear isometry.

  So \cref{lem:dc-ch8-2-1-jacobi-norm-transfer-manifold} applies and gives
  $|\tilde J(1)|=|J(1)|$, i.e. $|d(\exp_{\tilde p})_{i(v)}(i(Z))|=|d(\exp_p)_v(Z)|$. The
  chain rule along the semiconjugacy identifies $df_q(d(\exp_p)_v(Z))$ with
  $d(\exp_{\tilde p})_{i(v)}(i(Z))$ (\cref{lem:dc-ch8-2-1-mfderiv-bridge}), which turns that
  into $|df_q(w)|=|w|$ for every $w\in T_qM$; polarization
  (\cref{lem:dc-ch8-2-1-polarization}) upgrades the norm identity to the full inner product.
\end{proof}

\begin{lemma}[an orthonormal basis for the metric at a point]
  \label{lem:dc-ch8-2-1-metric-ortho-basis}
  \lean{Riemannian.Jacobi.exists_metricOrthonormalBasis}
  \leanok
  Let $g$ be a Riemannian metric on $M$ and $q\in M$. Then $T_qM$ has an orthonormal basis for
  $\langle\cdot,\cdot\rangle_q$.
\end{lemma}
\begin{proof}
  \leanok
\sketch
  $\langle\cdot,\cdot\rangle_q$ is a symmetric bilinear form on the finite-dimensional space
  $T_qM$, positive definite by the definition of a Riemannian metric. Diagonalize it and
  normalize.
\end{proof}

\begin{lemma}[Cartan transport and matched frames]
  \label{lem:dc-ch8-2-1-cartan-frame-data}
  \lean{Riemannian.Jacobi.cartanTransportZero, Riemannian.Jacobi.cartanSeed, Riemannian.Jacobi.metricInner_cartanSeed, Riemannian.Jacobi.cartanPhi, Riemannian.Jacobi.cartanPhi_zero, Riemannian.Jacobi.cartanPhi_velocity, Riemannian.Jacobi.cartanPhi_isLinear, Riemannian.Jacobi.metricInner_cartanPhi, Riemannian.Jacobi.exists_cartanFrameData}
  \uses{lem:dc-ch8-2-1-phi, lem:dc-ch8-2-1-transported-frame, lem:dc-ch8-2-1-parallel-transport-map, lem:dc-ch8-2-1-metric-ortho-basis, lem:dc-ch8-2-1-velocity-frame}
  \leanok
  Let $M,\tilde M$ be complete, $p\in M$, $\tilde p\in\tilde M$, let $i:T_pM\to T_{\tilde p}\tilde
  M$ be a linear isometry, let $v\in T_pM$, and let $[a',b']$ be a window with $a'<0<1<b'$. Write
  $\gamma=\gamma_v$ and $\tilde\gamma=\gamma_{i(v)}$.

  Define $j=\tilde P_0^{-1}\circ i\circ P_0$ and
  $\phi^{\mathrm{Cartan}}_t=\tilde P_t\circ j\circ P_t^{-1}$. Then for
  $t\in[a',b']$:
  \begin{itemize}
    \item $\phi^{\mathrm{Cartan}}_0=i$;
    \item $\phi^{\mathrm{Cartan}}_t(\gamma'(t))=\tilde\gamma'(t)$;
    \item $\phi^{\mathrm{Cartan}}_t$ is linear, and preserves the metric:
      $\langle\phi^{\mathrm{Cartan}}_tu,\phi^{\mathrm{Cartan}}_tw\rangle_{\tilde\gamma(t)}
      =\langle u,w\rangle_{\gamma(t)}$.
  \end{itemize}
  Moreover there exist parallel \emph{orthonormal} frames $e_1,\dots,e_n$ along $\gamma$ and
  $\tilde e_1,\dots,\tilde e_n$ along $\tilde\gamma$ on $[a',b']$ with
  $\phi^{\mathrm{Cartan}}_t(e_j(t))=\tilde e_j(t)$ for every $t\in[a',b']$.
\end{lemma}
\begin{proof}
  \leanok
  \uses{lem:dc-ch8-2-1-phi,lem:dc-ch8-2-1-transported-frame,
    lem:dc-ch8-2-1-metric-ortho-basis,lem:dc-ch8-2-1-velocity-frame}
\sketch
  Use the conjugated seed

  $$
  j=\tilde P_0^{-1}\circ i\circ P_0 ,
  $$

  Then $\phi_0=i$, and $j$ is a linear isometry because parallel transport and $i$
  preserve the metric.

  With $j$ in hand, seed the frames by an orthonormal basis at $\gamma(a')$
  (\cref{lem:dc-ch8-2-1-metric-ortho-basis}) and apply
  \cref{lem:dc-ch8-2-1-transported-frame}. The frame and velocity identities follow from
  uniqueness of parallel transport via \cref{lem:dc-ch8-2-1-phi}.
\end{proof}

\begin{lemma}[Cartan transfer with nonconjugacy discharged]
  \label{lem:dc-ch8-2-1-exp-norm-transfer-general-discharged}
  \lean{Riemannian.Jacobi.metricInner_mfderiv_eq_of_semiconjugacy_of_curvatureFormAt_of_isLocalDiffeomorphAt}
  \uses{lem:dc-ch8-2-1-exp-norm-transfer-general, lem:dc-ch8-2-1-cartan-frame-data, lem:dc-ch8-2-1-no-conjugate-normal}
  \leanok
  Same setting and same conclusion as \cref{lem:dc-ch8-2-1-exp-norm-transfer-general} --- $f$
  preserves the metric at $q=\exp_p(v)$ --- but with the frames, $\phi$, and \emph{both}
  no-conjugacy clauses discharged rather than assumed. In place of them it assumes only that
  $\exp_p$ and $\exp_{\tilde p}$ are local diffeomorphisms at $v$ and $i(v)$: the
  normal-neighborhood hypothesis. The remaining assumptions are the single-chart clauses
  and the curvature match for the Cartan transport of
  \cref{lem:dc-ch8-2-1-cartan-frame-data}.
\end{lemma}
\begin{proof}
  \leanok
  \uses{lem:dc-ch8-2-1-exp-norm-transfer-general,lem:dc-ch8-2-1-cartan-frame-data,
    lem:dc-ch8-2-1-no-conjugate-normal}
\sketch
  Obtain the frames and $\phi$ from \cref{lem:dc-ch8-2-1-cartan-frame-data} and feed the curvature
  hypothesis the produced $\phi$, which satisfies its three antecedents by construction. Obtain
  the two no-conjugacy clauses from \cref{lem:dc-ch8-2-1-no-conjugate-normal} applied on each
  side. Then apply \cref{lem:dc-ch8-2-1-exp-norm-transfer-general}.
\end{proof}

\begin{lemma}[from pointwise metric preservation to a local isometry on $V$]
  \label{lem:dc-ch8-2-1-local-isometry-packaging}
  \lean{Riemannian.DCIsLocalIsometryAt}
  \uses{lem:dc-ch8-2-1-exp-norm-transfer-general-discharged, lem:dc-ch8-2-1-single-chart, lem:dc-ch8-2-1-exp-diff-zero, lem:dc-ch8-2-1-no-conjugate-zero, lem:dc-ch8-2-1-window-transfer}
  \notready
  \notready
  The preceding result is pointwise at $q=\exp_p(v)$; this lemma packages it into a local
  isometry on the whole normal neighborhood. The remaining requirements are a single chart
  containing each relevant geodesic and the assembly of the pointwise differential identities.
  The base point satisfies $df_p=i$ by \cref{lem:dc-ch8-2-1-exp-diff-zero}.
\end{lemma}

\begin{lemma}[widening a single-chart time window]
  \label{lem:dc-ch8-2-1-window}
  \lean{Riemannian.Jacobi.exists_window_of_mem_chartAt_source}
  \leanok
  Let $\gamma:\mathbb R\to M$ be continuous and suppose $\gamma(t)$ lies in the source of a
  \emph{single} chart for every $t\in[0,1]$. Then it still does on some $[a',b']$ with $a'<0<1<b'$.

  This only widens a given single-chart hypothesis.
\end{lemma}
\begin{proof}
  \leanok
\sketch
  The preimage of the chart source is open and contains compact $[0,1]$, hence contains
  $(-\delta,1+\delta)$ for some $\delta>0$.
\end{proof}

\begin{lemma}[a geodesic of a normal neighborhood, confined to one chart]
  \label{lem:dc-ch8-2-1-single-chart}
  \uses{lem:dc-ch8-2-1-frame-chart-interface, lem:dc-ch8-2-1-jacobi-norm-transfer-general, lem:dc-ch8-2-1-jacobi-norm-transfer-manifold, lem:dc-ch8-2-1-chart-partition, lem:dc-ch8-2-1-radial-chart}
  \notready
  \notready
  The conclusion needed by the transfer is the single-chart condition on each geodesic over
  $[0,1]$ (or a finite chart-chain with compatible ODE data). The existence of such a chart
  uniformly for the normal neighborhood is the remaining obligation.
\end{lemma}

\begin{lemma}[the constant-curvature normal Jacobi field, any sign of $K_0$]
  \label{lem:dc-ch8-2-1-jacobi-const-solution}
  \lean{Riemannian.Jacobi.isJacobiFieldAlongOn_of_constantCurvature}
  \uses{def:dc-ch5-2-1}
  \leanok
  Let $M$ have constant sectional curvature $K_0$ (any sign), let $\gamma$ be a
  unit-speed geodesic, let $w$ be a \emph{parallel} field along $\gamma$ \emph{normal}
  to $\gamma'$, and let $h$ be any scalar solution of $h''+K_0 h=0$ (given by its
  derivative data $h'=Dh$, $Dh'=-K_0 h$). Then $J(t)=h(t)\,w(t)$ is a Jacobi field
  along $\gamma$, with covariant derivative $DJ(t)=Dh(t)\,w(t)$. This covers
  all curvature signs, with
  $h=\sin(t\sqrt{K_0})/\sqrt{K_0}$, $h=t$, $h=\sinh(t\sqrt{-K_0})/\sqrt{-K_0}$).
\end{lemma}
\begin{proof}
  \leanok
\sketch
  Since $w$ is parallel, $\tfrac{D}{dt}(h w)=h'\,w$ and
  $\tfrac{D}{dt}(Dh\,w)=Dh'\,w=-K_0 h\,w$; and by the constant-curvature model form
  $R(\gamma',w)\gamma'=K_0 w$ (using $|\gamma'|^2=1$ and
  $\langle w,\gamma'\rangle=0$) together with homogeneity in the middle slot,
  $R(\gamma',h w)\gamma'=h K_0\,w$. Hence
  $\tfrac{D^2}{dt^2}(h w)+R(\gamma',h w)\gamma'=-K_0 h\,w+K_0 h\,w=0$: the pair
  $(h w, Dh\,w)$ solves the Jacobi system. This generalizes the special case
  $K_0>0$, $h=\sin(t\sqrt{K_0})/\sqrt{K_0}$ used in
  \cref{ex:dc-ch5-3-3} to an arbitrary scalar solution.
\end{proof}

\begin{lemma}[the norm of a constant-curvature Jacobi field is manifold-independent]
  \label{lem:dc-ch8-2-1-jacobi-norm-transfer}
  \lean{Riemannian.Jacobi.metricInner_jacobiField_eq_of_constantCurvature_of_speedSq, Riemannian.Jacobi.metricInner_jacobiField_transfer_of_constantCurvature_of_speedSq}
  \uses{lem:dc-ch8-2-1-jacobi-const-solution, cor:dc-ch5-2-5, def:dc-ch5-2-1}
  \leanok
  Let $\gamma:[0,\ell]\to M$ be a geodesic of squared speed $c=|\gamma'|^2\neq0$ on a
  manifold of constant sectional curvature $K_0$, let $J$ be the Jacobi field with $J(0)=0$
  and covariant derivative $DJ$, and let $h$ solve $h''+K_0c\,h=0$, $h(0)=0$, $h'(0)=1$.
  Writing $Z=DJ(0)$ and $a=\langle Z,\gamma'(0)\rangle$,

  $$
  |J(t)|^2 = \frac{a^2}{c}\,t^2 + h(t)^2\Big(|Z|^2-\frac{a^2}{c}\Big).
  $$

  The right-hand side involves \emph{no} manifold data beyond $K_0$ and $c$ (through $h$),
  the time $t$, and the two scalar invariants $a=\langle Z,\gamma'(0)\rangle$ and $|Z|^2$.
  Consequently, Jacobi fields in any two spaces of curvature $K_0$ and equal speed $c$
  have equal norms when these two invariants agree. For $c=1$ this is
  $|J(t)|^2=t^2a^2+h(t)^2(|Z|^2-a^2)$.
\end{lemma}
\begin{proof}
  \leanok
\sketch
  Decompose $Z=\frac{a}{c}\,\gamma'(0)+Z_\perp$ with $Z_\perp=Z-\frac{a}{c}\gamma'(0)$; then
  $\langle Z_\perp,\gamma'(0)\rangle=a-\frac{a}{c}\cdot c=0$ and
  $|Z_\perp|^2=|Z|^2-\frac{a^2}{c}$. On a space of constant curvature $K_0$ a field normal to
  $\gamma'$ satisfies $R(\gamma',w)\gamma'=K_0|\gamma'|^2w=K_0c\,w$, so the normal Jacobi
  equation is the scalar equation $h''+K_0c\,h=0$. The tangential Jacobi field
  $\frac{a}{c}\,t\,\gamma'(t)$ and the normal Jacobi field
  $h(t)\,w(t)$, with $w$ the parallel transport of $Z_\perp$
  (\cref{lem:dc-ch8-2-1-jacobi-const-solution}), both vanish at $t=0$ and have covariant
  derivatives $\frac{a}{c}\gamma'(0)$ resp. $Z_\perp$ there; their sum has initial data
  $(0,Z)$, so by uniqueness of Jacobi fields with prescribed initial data it equals $J$.
  Parallel transport is an isometry preserving orthogonality, so
  $\langle\gamma'(t),w(t)\rangle=\langle\gamma'(0),Z_\perp\rangle=0$, $|\gamma'(t)|^2=c$
  (geodesics have constant speed) and $|w(t)|^2=|Z_\perp|^2=|Z|^2-\frac{a^2}{c}$; expanding
  the squared norm of the orthogonal combination
  $\big|\frac{a}{c}\,t\,\gamma'(t)+h(t)\,w(t)\big|^2$ gives
  $\frac{a^2}{c}t^2+h(t)^2(|Z|^2-\frac{a^2}{c})$, which proves the formula.
\end{proof}

\begin{lemma}[the differential of $\exp$ is norm-preserving between spaces of the same
  constant curvature]
  \label{lem:dc-ch8-2-1-exp-norm-transfer}
  \lean{Riemannian.Jacobi.chartMetricInner_expDifferential_eq_metricInner_jacobiField, Riemannian.Jacobi.chartMetricInner_expDifferential_transfer_of_constantCurvature_of_speedSq}
  \uses{lem:dc-ch8-2-1-jacobi-norm-transfer, cor:dc-ch5-2-5, def:dc-ch5-2-1}
  \leanok
  Let $M$ and $\tilde M$ be complete manifolds of the same constant sectional curvature
  $K_0$, let $p\in M$, $\tilde p\in\tilde M$, and let $v\in T_p M$, $\tilde v\in T_{\tilde p}\tilde M$
  be \emph{nonzero} vectors of the same length, $|v|_p^2=|\tilde v|_{\tilde p}^2=c$. Let
  $Z\in T_p M$ and $\tilde Z\in T_{\tilde p}\tilde M$ have matching scalar invariants,

  $$
  \langle \tilde Z,\tilde v\rangle_{\tilde p}=\langle Z,v\rangle_p,
  \qquad |\tilde Z|^2_{\tilde p}=|Z|^2_p .
  $$

  Then the two exponential differentials have equal norms,

  $$
  \big|(d\exp_{\tilde p})_{\tilde v}(\tilde Z)\big| = \big|(d\exp_p)_v(Z)\big| .
  $$

  If $i$ is a linear isometry with $\tilde v=i(v)$ and $\tilde Z=i(Z)$, these
  hypotheses hold automatically.
\end{lemma}
\begin{proof}
  \leanok
\sketch
  By \cref{cor:dc-ch5-2-5}, each differential is the value at $t=1$ of the
  Jacobi field with initial data $(0,Z)$ or $(0,\tilde Z)$. Constant speed gives the
  common squared speed $c$, so \cref{lem:dc-ch8-2-1-jacobi-norm-transfer} applies.
\end{proof}

\begin{lemma}[in constant curvature, no conjugate point where the scalar solution is nonzero]
  \label{lem:dc-ch8-2-1-no-conjugate}
  \lean{Riemannian.Jacobi.not_isConjugatePointAt_of_constantCurvature_of_speedSq, Riemannian.Jacobi.not_isConjugatePointAt_of_constantCurvature_of_lt_pi, Riemannian.Jacobi.exists_constCurvatureSol_ne_zero}
  \uses{lem:dc-ch8-2-1-jacobi-norm-transfer, lem:dc-ch8-2-1-jacobi-const-solution, def:dc-ch5-3-1}
  \leanok
  Let $\gamma:[0,\ell]\to M$ be a geodesic of squared speed $c\neq0$ on a manifold of constant
  sectional curvature $K_0$, and let $h$ solve $h''+K_0c\,h=0$, $h(0)=0$, $h'(0)=1$. If
  $h(\ell)\neq0$, then $\gamma(\ell)$ is \emph{not} conjugate to $\gamma(0)$ along $\gamma$.

  Only this direction is asserted. The converse --- that a zero of $h$ \emph{is} a conjugate
  point --- is not proved here; for $K_0>0$ it is witnessed at the first zero by
  \cref{ex:dc-ch5-3-3}, which exhibits $\gamma(\pi/\sqrt{K_0})$ as conjugate to $\gamma(0)$ in
  the unit-speed case (and additionally needs $n\geq2$).

  The criterion is sign-uniform. For $K_0\leq0$ the solution is $h(t)=t$ resp.
  $h(t)=\sinh(t\sqrt{-K_0c})/\sqrt{-K_0c}$, which never vanishes for $t>0$: there are no
  conjugate points at all, recovering the nonpositive-curvature statement of
  \cref{lem:dc-ch7-3-2}. For $K_0>0$ the solution is $h(t)=\sin(t\sqrt{K_0c})/\sqrt{K_0c}$,
  whose first positive zero is $\pi/\sqrt{K_0c}$, so the hypothesis holds throughout
  $0<\ell<\pi/\sqrt{K_0c}$ --- the interval that matters here, though $h$ is of course nonzero
  on further intervals beyond the first zero as well.

  Equivalently, in the $h$-free form: if

  $$
  K_0\,c\,\ell^2 < \pi^2,
  $$

  then $\gamma(\ell)$ is not conjugate to $\gamma(0)$. This single numerical condition covers
  all three signs at once --- it is vacuous for $K_0\leq0$, and for $K_0>0$ it reads
  $\ell\sqrt{K_0c}<\pi$.

  At parameter $1$ and $K_0>0$ this gives $|v|<\pi/\sqrt{K_0}$.
\end{lemma}
\begin{proof}
  \leanok
\sketch
  Let $J$ be a Jacobi field along $\gamma$ with $J(0)=0$ and $J(\ell)=0$; we show $J\equiv0$,
  which contradicts the nontriviality required of a conjugacy witness. Write $Z=DJ(0)$ and
  $a=\langle Z,\gamma'(0)\rangle$. By \cref{lem:dc-ch8-2-1-jacobi-norm-transfer},

  $$
  0=|J(\ell)|^2=\frac{a^2}{c}\,\ell^2+h(\ell)^2\Big(|Z|^2-\frac{a^2}{c}\Big).
  $$

  Both summands are nonnegative: $c>0$ because $c=|\gamma'(0)|^2$ and the metric is positive
  definite with $c\neq0$; and $|Z|^2-\frac{a^2}{c}\geq0$ by Cauchy--Schwarz, being the squared
  norm of the component of $Z$ orthogonal to $\gamma'(0)$. A sum of two nonnegative reals
  vanishes only if both do. The first gives $\frac{a^2}{c}\ell^2=0$, hence $a=0$ since
  $\ell>0$ and $c>0$. The second then reads $h(\ell)^2|Z|^2=0$, and $h(\ell)\neq0$ forces
  $|Z|^2=0$, so $Z=0$. Uniqueness then gives $J\equiv0$.

  For the $h$-free form, it remains to see that $K_0c\,\ell^2<\pi^2$ gives $h(\ell)\neq0$.
  Write $\kappa=K_0c$. If $\kappa<0$ then $h(t)=\sinh(t\sqrt{-\kappa})/\sqrt{-\kappa}$ and
  $\sinh$ is strictly increasing with $\sinh 0=0$, so $h(\ell)>0$; if $\kappa=0$ then
  $h(t)=t$ and $h(\ell)=\ell>0$; if $\kappa>0$ then $h(t)=\sin(t\sqrt\kappa)/\sqrt\kappa$, and
  $\kappa\ell^2<\pi^2$ with $\ell,\sqrt\kappa>0$ gives $0<\ell\sqrt\kappa<\pi$, whence
  $\sin(\ell\sqrt\kappa)>0$.
\end{proof}

\begin{lemma}[$\exp_p$ is a local diffeomorphism away from conjugate points]
  \label{lem:dc-ch8-2-1-exp-local-diffeo}
  \lean{Riemannian.Jacobi.isLocalDiffeomorphAt_expMapGlobal_of_not_conjugate}
  \uses{lem:dc-ch8-2-1-exp-equiv, prop:dc-ch5-3-5, lem:dc-ch7-3-2}
  \leanok
  Let $M$ be complete and $p\in M$. If the geodesic $\gamma_v$ from $p$ has no conjugate point
  of $p$ at parameter $1$, then $\exp_p$ is a $C^\infty$ local diffeomorphism at $v$; in
  particular it has a $C^\infty$ local inverse $\exp_p^{-1}$ defined near $\exp_p(v)$.

  Consequently $\exp_p^{-1}$ is smooth near $\exp_p(v)$.
\end{lemma}
\begin{proof}
  \leanok
\sketch
  By \cref{lem:dc-ch8-2-1-exp-equiv} the differential $d(\exp_p)_v$ is a continuous linear
  isomorphism: some chart $\zeta$ around $\exp_p(v)$ reads $\exp_p$ as a map with strict
  Fréchet derivative an isomorphism at $v$. That chart reading is $C^\infty$ on an open
  neighborhood of $v$, since $\exp_p$ is globally $C^\infty$ on a complete manifold, so the
  inverse function theorem makes the reading a local diffeomorphism at $v$. Composing with the
  chart inverse --- itself a local diffeomorphism --- recovers $\exp_p$ near $v$, which is
  therefore a local diffeomorphism at $v$ as well.
\end{proof}

\begin{lemma}[non-conjugacy from the normal neighborhood, in variable curvature]
  \label{lem:dc-ch8-2-1-no-conjugate-normal}
  \lean{Riemannian.Jacobi.not_isConjugatePointAt_globalGeodesic_of_injective_mfderiv, Riemannian.Jacobi.not_isConjugatePointAt_globalGeodesic_of_isLocalDiffeomorphAt, Riemannian.Jacobi.isLocalDiffeomorphAt_expMapGlobal_iff_not_isConjugatePointAt}
  \uses{lem:dc-ch8-2-1-exp-local-diffeo, lem:dc-ch8-2-1-mfderiv-bridge, prop:dc-ch5-3-5}
  \leanok
  Let $M$ be complete and $p\in M$, and let $v\in T_pM$. If $d(\exp_p)_v$ is injective --- in
  particular, if $\exp_p$ is a local diffeomorphism at $v$ --- then $1$ is \emph{not} a conjugate
  parameter along $\gamma_v$. No curvature hypothesis is needed.

  Together with \cref{lem:dc-ch8-2-1-exp-local-diffeo} this makes the two conditions
  \emph{equivalent}: $\exp_p$ is a $C^\infty$ local diffeomorphism at $v$ if and only if $1$ is
  not conjugate along $\gamma_v$.

  Thus the criterion applies in variable curvature whenever $\exp_p$ is locally invertible.
\end{lemma}
\begin{proof}
  \leanok
  \uses{lem:dc-ch8-2-1-mfderiv-bridge,prop:dc-ch5-3-5}
\sketch
  By \cref{prop:dc-ch5-3-5} it suffices to show that the chart reading $D$ of $d(\exp_p)_v$ is
  injective. \cref{lem:dc-ch8-2-1-mfderiv-bridge} factors it as

  $$
  D = C_{\exp_p(v)\to\zeta}\circ d(\exp_p)_v ,
  $$

  and the coordinate change $C_{\exp_p(v)\to\zeta}$ is injective, its cocycle inverse being a left
  inverse; so $D$ is a composite of two injective maps.

  A local diffeomorphism has an injective differential, giving the second assertion.
\end{proof}

\begin{lemma}[polarization: a linear map preserving lengths preserves the metric]
  \label{lem:dc-ch8-2-1-polarization}
  \lean{Riemannian.Jacobi.metricInner_polarization, Riemannian.Jacobi.metricInner_transfer_of_norm_transfer}
  \leanok
  Let $q\in M$, $\tilde q\in\tilde M$, and let $L:T_qM\to T_{\tilde q}\tilde M$ be additive. If
  $L$ preserves squared lengths, $|L(w)|^2_{\tilde q}=|w|^2_q$ for every $w\in T_qM$, then it
  preserves the metric itself:

  $$
  \langle L(u),L(v)\rangle_{\tilde q} = \langle u,v\rangle_q
  \qquad\text{for all } u,v\in T_qM .
  $$

  This is the step from \cref{lem:dc-ch8-2-1-exp-norm-transfer}, which delivers only the
  equality of \emph{lengths} $|df_q(u)|=|u|$, to the statement that $f$ is a local isometry,
  which is phrased with the full inner product. Additivity cannot be dropped: a
  length-preserving map that is not additive need not preserve the metric. In the application
  $L=df_q$ is linear, so the hypothesis is automatic.
\end{lemma}
\begin{proof}
  \leanok
\sketch
  A symmetric bilinear form over $\mathbb R$ is determined by its diagonal through the
  polarization identity: expanding $\langle u+v,u+v\rangle=\langle u,u\rangle+2\langle
  u,v\rangle+\langle v,v\rangle$ by bilinearity and symmetry gives

  $$
  \langle u,v\rangle=\tfrac12\big(\langle u+v,u+v\rangle-\langle u,u\rangle-\langle
  v,v\rangle\big).
  $$

  Apply this to $L(u),L(v)$ on $\tilde M$, use additivity to rewrite $L(u)+L(v)=L(u+v)$, then
  the length hypothesis at $u+v$, $u$ and $v$ to replace each of the three terms by its
  counterpart on $M$, and finally the same identity on $M$ read backwards:

  \langle L(u),L(v)\rangle=\tfrac12\big(|L(u+v)|^2-|L(u)|^2-|L(v)|^2\big)
  =\tfrac12\big(|u+v|^2-|u|^2-|v|^2\big)=\langle u,v\rangle. \qquad

\end{proof}

\begin{lemma}[$d(\exp_p)_v$ is an isomorphism, with the Jacobi identification]
  \label{lem:dc-ch8-2-1-exp-equiv}
  \lean{Riemannian.Jacobi.expDifferential_isEquiv_jacobi_of_not_conjugate}
  \uses{cor:dc-ch5-2-5, prop:dc-ch5-3-5}
  \leanok
  If the geodesic $\gamma_v$ from $p$ has no conjugate point of $p$ at parameter $1$, then
  $d(\exp_p)_v$ is a continuous linear isomorphism $T_pM\to T_{\exp_p(v)}M$, \emph{and} it still
  satisfies the identification of \cref{cor:dc-ch5-2-5}: it carries $\nabla J(0)$ to $J(1)$ for
  every Jacobi field $J$ along $\gamma_v$ with $J(0)=0$. Both facts come from the same strictly
  differentiable chart reading of $\exp_p$; the assembly of \cref{cor:dc-ch8-2-2} needs them
  simultaneously --- the isomorphism to invert $d\exp_p$ (so that every tangent vector at
  $q=\exp_p(v)$ is the endpoint value of a Jacobi field), and the identification to compute
  norms through \cref{lem:dc-ch8-2-1-exp-norm-transfer}.
\end{lemma}
\begin{proof}
  \leanok
\sketch
  \cref{prop:dc-ch5-3-5} makes the differential injective from the no-conjugate-point
  hypothesis; finite-dimensionality makes it an isomorphism. The Jacobi identification is
  \cref{cor:dc-ch5-2-5}.
\end{proof}

\begin{lemma}[the chart reading of a differential is the coordinate change of the intrinsic one]
  \label{lem:dc-ch8-2-1-mfderiv-bridge}
  \lean{Riemannian.Jacobi.chartReading_fderiv_apply_eq, Riemannian.Jacobi.chartMetricInner_chartReading_fderiv_eq_metricInner_mfderiv, Riemannian.Jacobi.chartMetricInner_expDifferential_eq_metricInner_mfderiv}
  \leanok
  Let $F:T_pM\to M$ be $C^\infty$, let $\zeta\in M$ be a chart base point with $F(v)$ in the
  chart source, and suppose the chart-$\zeta$ reading $w\mapsto\varphi_\zeta(F(w))$ has
  differential $D$ at $v$. Then $D$ is the intrinsic differential of $F$ followed by the
  coordinate change of its own target fibre:

  $$
  D(Z)=C_{F(v)\to\zeta}\big(dF_v(Z)\big)\qquad\text{for all }Z ,
  $$

  and consequently, for the chart Gram form,

  $$
  \langle D(Z),D(Z')\rangle_{\zeta,\varphi_\zeta(F(v))}=\langle dF_v(Z),dF_v(Z')\rangle_{F(v)} .
  $$

  This is the bridge between the chart-Gram form in which \cref{lem:dc-ch8-2-1-exp-norm-transfer}
  is stated and the intrinsic form that \cref{cor:dc-ch8-2-2} requires; applied to $F=\exp_p$ it
  turns every chart-level statement about $d(\exp_p)_v$ into an intrinsic one.
\end{lemma}
\begin{proof}
  \leanok
\sketch
  Apply the chain rule to $\varphi_\zeta\circ F$. Its differential is
  $C_{F(v)\to\zeta}\circ dF_v$, and uniqueness of the derivative gives the first
  display. Norm-faithfulness of chart readings gives the second.
\end{proof}

\begin{lemma}[$d(\exp_p)_0$ is the identity]
  \label{lem:dc-ch8-2-1-exp-diff-zero}
  \lean{Riemannian.Jacobi.mfderiv_expMapGlobal_zero_apply}
  \leanok
  For $M$ complete and $p\in M$, the differential of $\exp_p$ at the origin of $T_pM$ is the
  identity: $d(\exp_p)_0(v)=v$ for every $v\in T_pM$.
\end{lemma}
\begin{proof}
  \leanok
\sketch
  Fix $v$ and consider the ray $c:\mathbb R\to T_pM$, $c(t)=t\,v$, a linear map, so its own
  differential at $0$ is $c$ itself, and $c(1)=v$. Radial homogeneity of the exponential map gives
  $\exp_p(c(t))=\exp_p(t\,v)=\gamma_v(t)$, i.e. $\exp_p\circ\,c=\gamma_v$. Differentiating this
  identity at $t=0$ by the chain rule:
  $d(\exp_p)_0(c(1))=d(\gamma_v)_0(1)=\gamma_v'(0)=v$.
  Since $c(1)=v$, this is $d(\exp_p)_0(v)=v$.
\end{proof}

\begin{lemma}[the base point is never conjugate to itself at parameter $1$]
  \label{lem:dc-ch8-2-1-no-conjugate-zero}
  \lean{Riemannian.Jacobi.not_isConjugatePointAt_globalGeodesic_zero}
  \uses{lem:dc-ch8-2-1-mfderiv-bridge, lem:dc-ch8-2-1-exp-diff-zero, prop:dc-ch5-3-5}
  \leanok
  For $M$ complete and $p\in M$, the parameter $1$ is not conjugate along the constant geodesic
  $\gamma_0$. No curvature hypothesis is needed. Consequently, by
  \cref{lem:dc-ch8-2-1-exp-local-diffeo}, $\exp_p$ is a $C^\infty$ local diffeomorphism at
  $0\in T_pM$ on \emph{any} complete Riemannian manifold --- the normal-neighborhood fact.

\end{lemma}
\begin{proof}
  \leanok
\sketch
  By \cref{prop:dc-ch5-3-5} the parameter $1$ fails to be conjugate along $\gamma_v$ exactly when
  the chart reading $D$ of $d(\exp_p)_v$ is injective, so it suffices to prove that for $v=0$. By
  \cref{lem:dc-ch8-2-1-mfderiv-bridge}, $D(Z)=C_{\exp_p(0)\to\zeta}\big(d(\exp_p)_0(Z)\big)$, and
  $\exp_p(0)=p$ while $d(\exp_p)_0$ is the identity by \cref{lem:dc-ch8-2-1-exp-diff-zero}. Hence
  $D=C_{p\to\zeta}$, which is injective, since the cocycle identity
  $C_{\zeta\to p}\circ C_{p\to\zeta}=C_{p\to p}=\mathrm{id}$ exhibits a left inverse.
\end{proof}

\begin{lemma}[non-conjugacy along $\gamma_v$, and surjectivity of $d(\exp_p)_v$]
  \label{lem:dc-ch8-2-1-no-conjugate-ball}
  \lean{Riemannian.Jacobi.not_isConjugatePointAt_globalGeodesic_of_constantCurvature_of_lt_pi, Riemannian.Jacobi.surjective_mfderiv_expMapGlobal_of_not_conjugate}
  \uses{lem:dc-ch8-2-1-no-conjugate, lem:dc-ch8-2-1-exp-local-diffeo}
  \leanok
  Let $M$ have constant curvature $K_0$ and let $v\in T_pM$ with $v\neq 0$. If

  $$
  K_0\,\langle v,v\rangle_p<\pi^2 ,
  $$

  then $1$ is not conjugate along $\gamma_v$. The condition is vacuous for $K_0\leq 0$ and reads
  $|v|<\pi/\sqrt{K_0}$ for $K_0>0$; it is sharp there, by \cref{ex:dc-ch5-3-3}. Moreover, whenever
  $1$ is not conjugate along $\gamma_v$, the differential $d(\exp_p)_v$ is surjective onto
  $T_{\exp_p(v)}M$.

  Thus every tangent vector at $\exp_p(v)$ is $d(\exp_p)_v(Z)$ for some $Z$.
\end{lemma}
\begin{proof}
  \leanok
\sketch
  For the first claim, apply \cref{lem:dc-ch8-2-1-no-conjugate} along $\gamma_v$ at $\ell=1$: the
  geodesic $\gamma_v$ has constant speed $c=\langle v,v\rangle_p$, which is nonzero because the
  metric is positive definite and $v\neq 0$, and the numerical hypothesis
  $K_0\,c\,\ell^2<\pi^2$ is the assumption. For the second, \cref{lem:dc-ch8-2-1-exp-local-diffeo}
  makes $\exp_p$ a local diffeomorphism at $v$, and the differential of a local diffeomorphism is
  a linear isomorphism, in particular surjective.
\end{proof}

\begin{lemma}[the norm transfer, intrinsically]
  \label{lem:dc-ch8-2-1-exp-norm-transfer-mfderiv}
  \lean{Riemannian.Jacobi.metricInner_mfderiv_expMapGlobal_transfer_of_constantCurvature_of_speedSq}
  \uses{lem:dc-ch8-2-1-exp-norm-transfer, lem:dc-ch8-2-1-mfderiv-bridge, lem:dc-ch8-2-1-exp-equiv, lem:dc-ch8-2-1-jacobi-const-solution}
  \leanok
  Let $M$ and $\tilde M$ have the same constant curvature $K_0$, let $p\in M$, $\tilde p\in\tilde M$,
  and let $v,Z\in T_pM$, $\tilde v,\tilde Z\in T_{\tilde p}\tilde M$ with $v\neq 0$ and

  $$
  \langle v,v\rangle_p=\langle\tilde v,\tilde v\rangle_{\tilde p}=c,\qquad
  \langle \tilde Z,\tilde v\rangle_{\tilde p}=\langle Z,v\rangle_p,\qquad
  \langle \tilde Z,\tilde Z\rangle_{\tilde p}=\langle Z,Z\rangle_p .
  $$

  If $1$ is conjugate along neither $\gamma_v$ nor $\gamma_{\tilde v}$, then the two exponential
  differentials agree in length:

  $$
  \big|d(\exp_{\tilde p})_{\tilde v}(\tilde Z)\big|_{\exp_{\tilde p}(\tilde v)}
  =\big|d(\exp_p)_v(Z)\big|_{\exp_p(v)} .
  $$

  This is \cref{lem:dc-ch8-2-1-exp-norm-transfer} freed of its chart, and of the auxiliary scalar
  solution $h$.
\end{lemma}
\begin{proof}
  \leanok
\sketch
  \cref{lem:dc-ch8-2-1-exp-equiv} supplies, on each side, a chart around the endpoint in which the
  reading of the exponential map is differentiable with derivative $D$ (resp.\ $\tilde D$) still
  carrying $\nabla J(0)$ to $J(1)$. \cref{lem:dc-ch8-2-1-exp-norm-transfer} then equates the two
  chart-Gram norms of $D(Z)$ and $\tilde D(\tilde Z)$, its scalar hypothesis being discharged by
  \cref{lem:dc-ch8-2-1-jacobi-const-solution}, which produces a solution of $h''+(K_0c)h=0$ with
  $h(0)=0$, $h'(0)=1$ for every $K_0$ and $c$. Finally
  \cref{lem:dc-ch8-2-1-mfderiv-bridge} rewrites each chart-Gram norm as the intrinsic norm of the
  corresponding intrinsic differential.
\end{proof}

\begin{lemma}[metric preservation along a semiconjugacy]
  \label{lem:dc-ch8-2-2-semiconjugacy}
  \lean{Riemannian.Jacobi.mfderiv_apply_eq_of_semiconjugacy_zero, Riemannian.Jacobi.metricInner_mfderiv_eq_of_semiconjugacy, Riemannian.Jacobi.dcIsLocalIsometryAt_of_semiconjugacy, Riemannian.Jacobi.dcIsLocalIsometryAt_of_semiconjugacy_self}
  \uses{lem:dc-ch8-2-1-exp-norm-transfer-mfderiv, lem:dc-ch8-2-1-no-conjugate-ball, lem:dc-ch8-2-1-no-conjugate-zero, lem:dc-ch8-2-1-exp-diff-zero, lem:dc-ch8-2-1-polarization, lem:dc-ch8-2-1-exp-local-diffeo}
  \leanok
  Let $M$ and $\tilde M$ have the same constant curvature $K_0$, let $i:T_pM\to T_{\tilde p}\tilde M$
  be a linear isometry, and let $W\subset T_pM$ be an open neighborhood of $0$ on which
  $K_0\langle w,w\rangle_p<\pi^2$. Suppose $f:M\to\tilde M$ is differentiable at each $\exp_p(w)$,
  $w\in W$, and \emph{semiconjugates} $\exp_p$ to $\exp_{\tilde p}$ through $i$:

  $$
  f\big(\exp_p(w)\big)=\exp_{\tilde p}\big(i(w)\big)\qquad\text{for all }w\in W .
  $$

  Then $df_p=i$, and $f$ is a local isometry at $p$: there is a neighborhood $V$ of $p$ with
  $\langle u,u'\rangle_q=\langle df_q(u),df_q(u')\rangle_{f(q)}$ for all $q\in V$ and
  $u,u'\in T_qM$.

  Taking the semiconjugacy as the hypothesis, rather than $f=\exp_{\tilde p}\circ\,i\circ\exp_p^{-1}$,
  is what makes the differential of $\exp_p^{-1}$ unnecessary: $d(\exp_p)_v$ is surjective, so $df_q$
  is pinned down on all of $T_qM$ by its values on the image of $d(\exp_p)_v$.
\end{lemma}
\begin{proof}
  \leanok
\sketch
  Take $V=\exp_p(W')$, where $W'\subset W$ is the part of $W$ on which $\exp_p$ is inverted by the
  local inverse supplied by \cref{lem:dc-ch8-2-1-exp-local-diffeo} at $0$, available by
  \cref{lem:dc-ch8-2-1-no-conjugate-zero}. This $V$ is a neighborhood of $p$, and each $q\in V$ is
  $\exp_p(w)$ for a unique $w\in W'$.

  Suppose first $w\neq0$. Since $i$ is a linear isometry, $\langle i(w),i(w)\rangle_{\tilde p}
  =\langle w,w\rangle_p$, so the numerical hypothesis holds on both sides and by
  \cref{lem:dc-ch8-2-1-no-conjugate-ball} the parameter $1$ is conjugate along neither $\gamma_w$
  nor $\gamma_{i(w)}$; in particular $d(\exp_p)_w$ is surjective. Differentiating the semiconjugacy
  at $w$ by the chain rule --- both sides are differentiable, and the derivative of a map is unique
  --- gives

  $$
  df_q\big(d(\exp_p)_w(Z)\big)=d(\exp_{\tilde p})_{i(w)}\big(i(Z)\big)\qquad\text{for all }Z .
  $$

  The invariants $\langle i(Z),i(w)\rangle=\langle Z,w\rangle$ and $|i(Z)|^2=|Z|^2$ are preserved
  because $i$ is a linear isometry, so \cref{lem:dc-ch8-2-1-exp-norm-transfer-mfderiv} equates the
  lengths of the two sides. As $Z$ ranges over $T_pM$ the vector $d(\exp_p)_w(Z)$ ranges over all of
  $T_qM$, so $df_q$ preserves lengths; being linear, it preserves the metric, by
  \cref{lem:dc-ch8-2-1-polarization}.

  If instead $w=0$ then $q=\exp_p(0)=p$. The same chain rule at $0$, together with
  $d(\exp_p)_0=d(\exp_{\tilde p})_0=\mathrm{id}$ (\cref{lem:dc-ch8-2-1-exp-diff-zero}) and
  $i(0)=0$, gives $df_p=i$; and $f(p)=\exp_{\tilde p}(i(0))=\tilde p$. The claim at $q=p$ is then
  exactly the hypothesis that $i$ is a linear isometry.
\end{proof}

\begin{theorem}[Cartan local isometry theorem]
  \label{thm:dc-ch8-2-1}
  \uses{cor:dc-ch5-2-5, lem:dc-ch8-2-1-jacobi-transfer, lem:dc-ch8-2-1-parallel-ortho-frame, lem:dc-ch8-2-1-transported-frame, lem:dc-ch8-2-1-hmatch-transfer, lem:dc-ch8-2-1-curvature-bridge, lem:dc-ch8-2-1-frame-coeff-bridge, lem:dc-ch8-2-1-jacobi-norm-transfer-general, lem:dc-ch8-2-1-frame-chart-interface, lem:dc-ch8-2-1-jacobi-norm-transfer-manifold, lem:dc-ch8-2-1-single-chart, lem:dc-ch8-2-1-exp-norm-transfer-general, lem:dc-ch8-2-1-phi, lem:dc-ch8-2-1-parallel-transport-map, lem:dc-ch8-2-1-exp-norm-transfer, lem:dc-ch8-2-1-exp-equiv, lem:dc-ch8-2-1-exp-local-diffeo, lem:dc-ch8-2-1-polarization, lem:dc-ch8-2-1-no-conjugate, lem:dc-ch8-2-1-mfderiv-bridge, lem:dc-ch8-2-1-exp-diff-zero, lem:dc-ch8-2-1-no-conjugate-zero, lem:dc-ch8-2-1-no-conjugate-ball, lem:dc-ch8-2-1-exp-norm-transfer-mfderiv, lem:dc-ch8-2-1-jacobi-reversal, lem:dc-ch8-2-1-jacobi-backward-uniqueness, lem:dc-ch8-2-1-jacobi-interior-data, lem:dc-ch8-2-1-local-isometry-packaging}
  \notready
  \dcref{ch8:2.1}
  \notready
  With the notation above, if for all $q\in V$ and all $x,y,u,v\in T_q(M)$ we have


  $$
  \langle R(x,y)u,v\rangle=\big\langle\tilde{R}(\phi_t(x),\phi_t(y))\phi_t(u),\phi_t(v)\big\rangle,
  $$


  then $f:V\to f(V)\subset\tilde{M}$ is a local isometry and $df_p=i$.
\end{theorem}
\begin{proof}
\sketch
  Let $q\in V$ and let $\gamma:[0,\ell]\to M$ be a normalized geodesic with
  $\gamma(0)=p$, $\gamma(\ell)=q$. Let $v\in T_q(M)$ and let $J$ be a Jacobi field
  along $\gamma$ with $J(0)=0$ and $J(\ell)=v$. Let $e_1,\dots,e_n=\gamma'(0)$ be an
  orthonormal basis of $T_p(M)$ and let $e_i(t)$ be the parallel transport of $e_i$
  along $\gamma$. Writing $J(t)=\sum_i y_i(t)e_i(t)$, the Jacobi equation gives


  $$
  y_j''+\sum_i\langle R(e_n,e_i)e_n,e_j\rangle y_i=0,\qquad j=1,\dots,n.
  $$


  Let $\tilde{\gamma}:[0,\ell]\to\tilde{M}$ be the normalized geodesic with
  $\tilde{\gamma}(0)=\tilde{p}$, $\tilde{\gamma}'(0)=i(\gamma'(0))$, and set
  $\tilde{J}(t)=\phi_t(J(t))$, $\tilde{e}_j(t)=\phi_t(e_j(t))$. By linearity of
  $\phi_t$, $\tilde{J}(t)=\sum_i y_i(t)\tilde{e}_i(t)$. Since by hypothesis
  $\langle R(e_n,e_i)e_n,e_j\rangle=\langle\tilde{R}(\tilde{e}_n,\tilde{e}_i)\tilde{e}_n,\tilde{e}_j\rangle$,
  $\tilde{J}$ satisfies the same equation, hence $\tilde{J}$ is a Jacobi field along
  $\tilde{\gamma}$ with $\tilde{J}(0)=0$. Since parallel transport is an isometry,
  $|\tilde{J}(\ell)|=|J(\ell)|$. It remains to show that $\tilde{J}(\ell)=df_q(v)=df_q(J(\ell))$.
  Since $\tilde{J}(t)=\phi_t(J(t))$, $\tilde{J}'(0)=i(J'(0))$. Since $J$ and
  $\tilde{J}$ are Jacobi fields vanishing at $t=0$, by \cref{cor:dc-ch5-2-5},


  $$
  J(t)=(d\exp_p)_{t\gamma'(0)}(tJ'(0)),\qquad\tilde{J}(t)=(d\exp_{\tilde{p}})_{t\tilde{\gamma}'(0)}(t\tilde{J}'(0)).
  $$


  Therefore


  $$
  \tilde{J}(\ell)=(d\exp_{\tilde{p}})_{\ell\tilde{\gamma}'(0)}\circ i\circ\big((d\exp_p)_{\ell\gamma'(0)}\big)^{-1}(J(\ell))=df_q(J(\ell)),
  $$


  which proves the assertion.

  If both exponential maps are global diffeomorphisms, the same argument makes $f$
  a global isometry.
\end{proof}

\begin{corollary}[Local equivalence of equal-curvature spaces]
  \label{cor:dc-ch8-2-2}
  \lean{Riemannian.Jacobi.exists_dcIsLocalIsometryAt_of_constantCurvature_of_orthonormal_bases}
  \uses{lem:dc-ch8-2-1-jacobi-transfer-const, lem:dc-ch8-2-1-velocity-frame, lem:dc-ch8-2-1-jacobi-norm-transfer, lem:dc-ch8-2-1-jacobi-const-solution, lem:dc-ch8-2-1-exp-norm-transfer, lem:dc-ch8-2-1-exp-equiv, lem:dc-ch8-2-1-exp-local-diffeo, lem:dc-ch8-2-1-polarization, lem:dc-ch8-2-1-no-conjugate, lem:dc-ch8-2-1-mfderiv-bridge, lem:dc-ch8-2-1-exp-diff-zero, lem:dc-ch8-2-1-no-conjugate-zero, lem:dc-ch8-2-1-no-conjugate-ball, lem:dc-ch8-2-1-exp-norm-transfer-mfderiv, lem:dc-ch8-2-2-semiconjugacy}
  \leanok
  \dcref{ch8:2.2}
  Let $M$ and $\tilde{M}$ be spaces with the same constant curvature and the same
  dimension $n$. Let $p\in M$ and $\tilde{p}\in\tilde{M}$. Choose arbitrary
  orthonormal bases $\{e_j\}$ of $T_p(M)$ and $\{\tilde{e}_j\}$ of $T_{\tilde{p}}(\tilde{M})$,
  $j=1,\dots,n$. Then there exist a neighborhood $V\subset M$ of $p$, a neighborhood
  $\tilde{V}\subset\tilde{M}$ of $\tilde{p}$, and an isometry $f:V\to\tilde{V}$ such
  that $df_p(e_j)=\tilde{e}_j$.
\end{corollary}
\begin{proof}
  \leanok
\sketch
  Let $i(e_j)=\tilde e_j$ and define
  $f=\exp_{\tilde p}\circ i\circ\exp_p^{-1}$ on a sufficiently small normal
  neighborhood. The constant-curvature identity gives the curvature match; equivalently,
  \cref{lem:dc-ch8-2-2-semiconjugacy} applies on
  $K_0\langle w,w\rangle_p<\pi^2$. It yields $df_p=i$ and local metric preservation.
  Restricting the resulting local diffeomorphism gives $f:V\to\tilde V$.
\end{proof}

\begin{corollary}[Local frame transitivity]
  \label{cor:dc-ch8-2-3}
  \lean{Riemannian.Jacobi.exists_dcIsLocalIsometryAt_of_constantCurvature_self_of_orthonormal_bases}
  \uses{lem:dc-ch8-2-1-jacobi-transfer-const, lem:dc-ch8-2-1-velocity-frame, lem:dc-ch8-2-1-jacobi-norm-transfer, lem:dc-ch8-2-1-jacobi-const-solution, lem:dc-ch8-2-1-exp-norm-transfer, lem:dc-ch8-2-1-exp-equiv, lem:dc-ch8-2-1-exp-local-diffeo, lem:dc-ch8-2-1-polarization, lem:dc-ch8-2-1-no-conjugate, lem:dc-ch8-2-1-mfderiv-bridge, lem:dc-ch8-2-1-exp-diff-zero, lem:dc-ch8-2-1-no-conjugate-zero, lem:dc-ch8-2-1-no-conjugate-ball, lem:dc-ch8-2-1-exp-norm-transfer-mfderiv, lem:dc-ch8-2-2-semiconjugacy}
  \leanok
  \dcref{ch8:2.3}
  Let $M$ be a space of constant curvature and let $p$ and $q$ be any two points of
  $M$. Let $\{e_j\}$ and $\{f_j\}$ be arbitrary orthonormal bases of $T_p(M)$ and
  $T_q(M)$, respectively. Then there exist neighborhoods $U$ of $p$ and $V$ of $q$,
  and an isometry $g:U\to V$ such that $dg_p(e_j)=f_j$.

\end{corollary}
\begin{proof}
  \leanok
  \uses{cor:dc-ch8-2-2}
\sketch
  Apply \cref{cor:dc-ch8-2-2} with $\tilde M=M$ and $\tilde p=q$, taking for $i$ the linear
  isometry $T_pM\to T_qM$ carrying $e_j$ to $f_j$.
\end{proof}

\section{Hyperbolic space}

\begin{definition}[hyperbolic space $H^n$ as a Riemannian manifold]
  \label{def:dc-ch8-3-hyperbolic-space}
  \lean{Riemannian.opensConformalMetric, Riemannian.Hyperbolic.hyperbolicMetric, Riemannian.Hyperbolic.hyperbolicMetric_apply, Riemannian.Hyperbolic.hyperbolicMetric_gram}
  The \emph{hyperbolic space} $H^n$ is the open half-space
  $H^n=\{x\in\mathbb R^n : x_n>0\}$ carrying the metric $g_{ij}=\delta_{ij}/x_n^2$
  of $(1)$, that is, the conformal rescaling of the ambient Euclidean metric by
  the positive smooth factor $1/x_n^2$: for $u,v\in T_pH^n$,
  $\langle u,v\rangle_g=\tfrac{1}{x_n^2}\langle u,v\rangle$.
\end{definition}

\begin{lemma}[Christoffel symbols and coordinate curvature of a conformal metric]
  \label{lem:dc-ch8-3-conformal-curvature}
  \lean{Riemannian.Hyperbolic.christoffelFromMetric_eq, Riemannian.Hyperbolic.Rcoeff_diag, Riemannian.Hyperbolic.hyperbolic_sectionalCurvature}
  \leanok
  For the conformal metric $g_{ij}=\delta_{ij}/F^2$ ($f=\log F$), the Christoffel
  symbols computed from the Koszul formula
  $\Gamma^k_{ij}=\tfrac12\sum_\ell g^{k\ell}(\partial_i g_{j\ell}+\partial_j g_{i\ell}-\partial_\ell g_{ij})$
  equal $\Gamma^k_{ij}=-\delta_{jk}f_i-\delta_{ki}f_j+\delta_{ij}f_k$, and the
  associated coordinate curvature coefficients satisfy, for $i\neq j$,
  $$F^2R_{ijij}=-\sum_\ell f_\ell^2+f_i^2+f_j^2+f_{ii}+f_{jj}.$$
  Specializing to $F=x_n$ (so $f=\log x_n$, $f_a=\delta_{an}/x_n$,
  $f_{ab}=-\delta_{an}\delta_{bn}/x_n^2$) gives, for every \emph{coordinate}
  $2$-plane spanned by $\partial/\partial x_i,\partial/\partial x_j$ with $i\neq j$,
  the coordinate-frame sectional curvature $K_{ij}=R_{ijij}F^4=-1$. (This is the
  coordinate identity; extending it to all $2$-planes and to the intrinsic
  curvature of the manifold is \cref{prop:dc-ch8-3-const-curv}.)
\end{lemma}
\begin{proof}
  \leanok
\sketch
  The Koszul identity is a direct computation from
  $g^{k\ell}=\delta_{k\ell}F^2$ and $\partial_m g_{ab}=-2\delta_{ab}f_m/F^2$.
  The curvature coefficient is
  $R^s_{ijk}=\partial_j\Gamma^s_{ik}-\partial_i\Gamma^s_{jk}+\sum_\ell(\Gamma^\ell_{ik}\Gamma^s_{j\ell}-\Gamma^\ell_{jk}\Gamma^s_{i\ell})$;
  substituting the closed form for $\Gamma$ and contracting the Kronecker deltas
  yields $R^j_{iji}=-\sum_\ell f_\ell^2+f_i^2+f_j^2+f_{ii}+f_{jj}$, which is
  $F^2R_{ijij}$ since $g_{jj}=1/F^2$. For $F=x_n$ one has
  $(f_a)^2+f_{aa}=0$ at every index, so the bracketed sum collapses to
  $-\sum_\ell f_\ell^2=-1/x_n^2$, and $K_{ij}=R_{ijij}F^4=(F^2R_{ijij})F^2=-1$.
\end{proof}

\begin{lemma}[the Riemann curvature tensor in chart coordinates]
  \label{lem:dc-ch8-3-chart-curvature}
  \lean{Riemannian.leviCivita_curvature_chartFrame_expansion}
  \leanok
  Let $g$ be a Riemannian metric on $M$ and $\nabla$ its Levi-Civita connection.
  In the coordinate frame $\partial_i$ of a chart at a point $p$, the intrinsic
  curvature operator is computed from the chart Christoffel symbols
  $\Gamma^k_{ij}$ by the classical coordinate formula

  $$
  R(\partial_i,\partial_j)\partial_k
    = \sum_\ell\Big(\partial_j\Gamma^\ell_{ik}-\partial_i\Gamma^\ell_{jk}
      +\sum_s(\Gamma^s_{ik}\Gamma^\ell_{js}-\Gamma^s_{jk}\Gamma^\ell_{is})\Big)\partial_\ell,
  $$

  where $\partial_j\Gamma^\ell_{ik}$ is the chart partial derivative of the
  coordinate Christoffel function. With the convention
  $R(X,Y)Z=\nabla_Y\nabla_X Z-\nabla_X\nabla_Y Z+\nabla_{[X,Y]}Z$ the bracketed
  coefficient is exactly the $R^\ell_{ijk}$ of $(2)$.
\end{lemma}
\begin{proof}
  \leanok
\sketch
  In a coordinate chart the frame vectors have vanishing pairwise Lie brackets,
  $[\partial_i,\partial_j]=0$, so the bracket term
  $\nabla_{[\partial_i,\partial_j]}\partial_k$ drops out. The field
  $\nabla_{\partial_j}\partial_k$ is, throughout the chart, the frame combination
  $\sum_m\Gamma^m_{jk}\partial_m$ (formula $(10)$ of Chap.~2, the Koszul collapse
  on the coordinate frame, valid at every chart point). Differentiating this
  along $\partial_i$ by the Leibniz rule for $\nabla$ gives
  $\nabla_{\partial_i}\nabla_{\partial_j}\partial_k
   =\sum_\ell\big(\partial_i\Gamma^\ell_{jk}
   +\sum_s\Gamma^s_{jk}\Gamma^\ell_{is}\big)\partial_\ell$, using
  $\nabla_{\partial_i}\partial_m=\sum_\ell\Gamma^\ell_{im}\partial_\ell$ for the
  inner derivative. Subtracting the term with $i$ and $j$ interchanged yields the
  stated coefficients.
\end{proof}

\begin{lemma}[hyperbolic chart Gram and inverse-Gram in the abstract chart frame]
  \label{lem:dc-ch8-3-hyperbolic-chart-gram}
  \lean{Riemannian.Hyperbolic.hyperbolic_chartGramMatrix, Riemannian.Hyperbolic.hyperbolic_chartGramMatrix_eq, Riemannian.Hyperbolic.hyperbolic_chartInvGramMatrix, Riemannian.Hyperbolic.finBasisGramMatrix_det_isUnit}
  \uses{def:dc-ch8-3-hyperbolic-space}
  \leanok
  Because $H^n=\{x_n>0\}$ is an open subset of $\mathbb R^n$, its tangent bundle is
  canonically trivial and the abstract chart frame vector
  $\partial_i=$ \emph{finBasis}$_i$ (the frame underlying \cref{lem:dc-ch8-3-chart-curvature})
  is literally the constant basis vector. Hence the chart Gram matrix of the
  hyperbolic metric is the conformal rescaling
  $G_{ij}(x)=x_n^{-2}\,B_{ij}$ of the constant Gram matrix
  $B_{ij}=\langle\text{finBasis}_i,\text{finBasis}_j\rangle$ of the abstract basis,
  and its inverse is $G^{ij}(x)=x_n^{2}\,B^{ij}$ (with $B$ invertible, being the
  Gram matrix of a basis). Since the abstract basis need not be
  $\langle\,,\rangle$-orthonormal, $B$ is a \emph{general} constant symmetric
  positive-definite matrix, not $\delta$: $\delta_{ij}$ is
  replaced by $B_{ij}$ in the abstract chart frame, and the conformal
  computation goes through with $\delta\rightsquigarrow B$, the extra factors
  collapsing by the inverse-Gram identity $\sum_{ij}B^{ij}c_ic_j=1$ (for
  $c_i=\langle e_n,\text{finBasis}_i\rangle$).
\end{lemma}
\begin{proof}
  \leanok
\sketch
  The trivialisation of the tangent bundle of an open subset of $\mathbb R^n$ is
  the identity on fibres, so $\partial_i(x)=$ finBasis$_i$ and
  $G_{ij}(x)=\langle\partial_i,\partial_j\rangle_g=x_n^{-2}\langle\text{finBasis}_i,\text{finBasis}_j\rangle=x_n^{-2}B_{ij}$.
  The Gram matrix at the basepoint is positive definite, so $\det B$ is a unit and
  the inverse of the scalar multiple $x_n^{-2}B$ is $x_n^{2}B^{-1}$.
\end{proof}

\begin{lemma}[hyperbolic chart Christoffel symbols and their derivatives, closed form]
  \label{lem:dc-ch8-3-hyperbolic-christoffel}
  \lean{Riemannian.Hyperbolic.hyperbolic_chartChristoffel, Riemannian.Hyperbolic.hyperbolic_partialDeriv_chartChristoffel, Riemannian.Hyperbolic.hyperbolic_partialDeriv_chartGramOnE}
  \uses{lem:dc-ch8-3-hyperbolic-chart-gram}
  \leanok
  In the abstract chart frame the Christoffel symbols of the hyperbolic metric are,
  in closed form,
  $$\Gamma^k_{ij}(y)=-f_i\delta^k_j-f_j\delta^k_i+d^k B_{ij},\qquad
    f_i=\frac{c_i}{y_n},\ c_i=(\text{finBasis}_i)_n,\ d^k=\frac{\sum_\ell B^{k\ell}c_\ell}{y_n},$$
  the constant abstract-basis Gram $B$ (\cref{lem:dc-ch8-3-hyperbolic-chart-gram})
  replacing $\delta$ in the abstract frame.
  Because every term carries the single factor $1/y_n$, one may write
  $\Gamma^k_{ij}(y)=K^k_{ij}\,(y_n)^{-1}$ with $K^k_{ij}$ the ($y$-independent)
  numerator, so the abstract-chart-frame partial derivatives are
  $$\partial_m\Gamma^k_{ij}(y)=K^k_{ij}\cdot\bigl(-c_m\,(y_n)^{-2}\bigr).$$
  This is exactly the $\Gamma$ and $\partial\Gamma$ data consumed by the chart-frame
  curvature expansion \cref{lem:dc-ch8-3-chart-curvature}.
\end{lemma}
\begin{proof}
  \leanok
\sketch
  The chart Gram is $G_{ij}(y)=(y_n^2)^{-1}B_{ij}$ and its inverse
  $G^{k\ell}(y)=y_n^2 B^{k\ell}$ (\cref{lem:dc-ch8-3-hyperbolic-chart-gram}). The
  distinguished coordinate $y_n=y_e$ has derivative
  $\partial_m y_n=(\text{finBasis}_m)_n=c_m$ along $\text{finBasis}_m$, so
  $\partial_m\bigl((y_n^2)^{-1}\bigr)=-2\,y_n^{-3}c_m$ and hence
  $\partial_m G_{ij}=B_{ij}\cdot(-2\,y_n^{-3}c_m)$. Substituting into the Koszul
  formula
  $\Gamma^k_{ij}=\tfrac12\sum_\ell G^{k\ell}(\partial_iG_{j\ell}+\partial_jG_{i\ell}-\partial_\ell G_{ij})$
  and contracting $\sum_\ell B^{k\ell}B_{\ell j}=\delta^k_j$,
  $\sum_\ell B^{k\ell}B_{\ell i}=\delta^k_i$ (from $B^{-1}B=1$) collapses the factors
  $y_n^2\cdot y_n^{-3}=y_n^{-1}$, giving the closed form. Since only the $1/y_n$ factor
  depends on $y$, $\partial_m\Gamma^k_{ij}=K^k_{ij}\,\partial_m\bigl(y_n^{-1}\bigr)
  =K^k_{ij}\,(-c_m y_n^{-2})$.
\end{proof}

\begin{lemma}[inverse-Gram/coordinate identity $\sum_{ij}B^{ij}c_ic_j=1$]
  \label{lem:dc-ch8-3-inverse-gram-coord}
  \lean{Riemannian.Hyperbolic.finBasis_inv_gram_coord}
  \uses{lem:dc-ch8-3-hyperbolic-chart-gram}
  \leanok
  For the constant abstract-basis Gram matrix
  $B_{ij}=\langle\text{finBasis}_i,\text{finBasis}_j\rangle$ with inverse
  $B^{ij}=(B^{-1})_{ij}$, and $c_i=(\text{finBasis}_i)_n$ the $n$-th coordinate of
  the $i$-th basis vector,
  $$\sum_{i,j} B^{ij}\,c_i\,c_j = 1.$$
  Equivalently $c_i=\langle\text{finBasis}_i, w_n\rangle$ for the standard unit
  vector $w_n=\text{basisFun}_n$, and the $B^{-1}$-contraction reconstructs
  $\langle w_n, w_n\rangle=1$.
\end{lemma}
\begin{proof}
  \leanok
\sketch
  Write $w_n=\sum_i a_i\,\text{finBasis}_i$ in the abstract basis. Then
  $c_i=\langle\text{finBasis}_i,w_n\rangle=\sum_j
  B_{ij}a_j$, so $\sum_j B^{ij}c_j=(B^{-1}B a)_i=a_i$ and
  $\sum_{ij}B^{ij}c_ic_j=\sum_i a_i c_i=\langle w_n,w_n\rangle=1$.
\end{proof}

\begin{lemma}[the pointwise $(0,4)$ curvature form and its chart-frame value]
  \label{lem:dc-ch8-3-pointwise-curvature-form}
  \lean{Riemannian.AffineConnection.curvatureFormAt, Riemannian.AffineConnection.isAlgCurvatureForm_curvatureFormAt, Riemannian.leviCivita_curvatureFormAt_chartFrame}
  \uses{lem:dc-ch8-3-chart-curvature}
  \leanok
  The curvature operator $R(X,Y)Z$ is a tensor: its value at $p$ depends only on
  $X(p),Y(p),Z(p)$. Lowering by the
  metric gives the pointwise $(0,4)$ form $\langle R(x,y)z,t\rangle_g$ on
  $T_pM$, which for the Levi-Civita connection is an \emph{algebraic curvature
  form} (multilinear, antisymmetric in each pair, and satisfying first Bianchi).
  In the chart frame $\partial_a=\text{finBasis}_a$ its value is the
  coordinate curvature coefficient contracted with the chart Gram,
  $$\langle R(\partial_i,\partial_j)\partial_k,\partial_\ell\rangle_g
    = \sum_m R^m_{ijk}\,G_{m\ell},\qquad
    R^m_{ijk}=\partial_j\Gamma^m_{ik}-\partial_i\Gamma^m_{jk}
      +\sum_s(\Gamma^s_{ik}\Gamma^m_{js}-\Gamma^s_{jk}\Gamma^m_{is}),$$
  the metric-lowered avatar of \cref{lem:dc-ch8-3-chart-curvature}. This bridges
  the \emph{intrinsic} manifold curvature to the coordinate Christoffel data, and
  is the input to the all-planes identification \cref{cor:dc-ch4-3-5}: it makes the
  chart coefficients the components of an algebraic curvature form on $T_pM$.
\end{lemma}
\begin{proof}
  \leanok
\sketch
  Tensoriality is the pointwise reading of the $\mathcal D(M)$-multilinearity of
  $R$: a field vanishing at $p$ decomposes as $\sum_i f_i W_i+\tau$ with
  $f_i(p)=0$ and $\tau$ vanishing near $p$, and homogeneity of $R$ in that slot
  kills both parts, so $R(\cdot)$ against a field vanishing at $p$ vanishes at
  $p$; hence $R(X,Y)Z|_p$ depends only on the values at $p$. The
  algebraic-curvature-form axioms transport slotwise from the field-level
  symmetries (Prop.~2.5, via the Levi-Civita property). The chart-frame value
  pairs the frame expansion of $R(\partial_i,\partial_j)\partial_k$
  (\cref{lem:dc-ch8-3-chart-curvature}) against $\partial_\ell$, distributing the
  metric over the finite sum.
\end{proof}

\begin{proposition}[$H^n$ has constant sectional curvature $-1$]
  \label{prop:dc-ch8-3-const-curv}
  \lean{Riemannian.Hyperbolic.hyperbolic_isConstantCurvature, Riemannian.Hyperbolic.hyperbolicMetric_sectionalCurvature_eq_neg_one, Riemannian.Hyperbolic.hyperbolic_curvatureFormAt_eq, Riemannian.Hyperbolic.hyperbolic_curvatureFormAt_frame}
  \uses{lem:dc-ch8-3-conformal-curvature, lem:dc-ch8-3-chart-curvature, lem:dc-ch8-3-pointwise-curvature-form, lem:dc-ch8-3-hyperbolic-chart-gram, lem:dc-ch8-3-hyperbolic-christoffel, lem:dc-ch8-3-inverse-gram-coord, cor:dc-ch4-3-5, def:dc-ch8-3-hyperbolic-space}
  \leanok
  The hyperbolic space $H^n$, with the metric $g_{ij}=\delta_{ij}/x_n^2$, has
  constant sectional curvature equal to $-1$: for every point and every
  $2$-plane $\sigma\subset T_pH^n$, $K(p,\sigma)=-1$.
\end{proposition}
\begin{proof}
  \leanok
\sketch
  Work in the abstract chart frame $\partial_a=\text{finBasis}_a$ of $T_pH^n$. By
  \cref{lem:dc-ch8-3-pointwise-curvature-form} the pointwise $(0,4)$ curvature form
  of the Levi-Civita connection is, in this frame,
  $\langle R(\partial_i,\partial_j)\partial_k,\partial_\ell\rangle_g=\sum_m
  R^m_{ijk}G_{m\ell}$, with $R^m_{ijk}=\partial_j\Gamma^m_{ik}-\partial_i\Gamma^m_{jk}
  +\sum_s(\Gamma^s_{ik}\Gamma^m_{js}-\Gamma^s_{jk}\Gamma^m_{is})$ and $G_{m\ell}$
  the chart Gram matrix. Insert the closed forms
  \cref{lem:dc-ch8-3-hyperbolic-christoffel}, $\Gamma^k_{ij}(y)=y_n^{-1}K^k_{ij}$
  and $\partial_m\Gamma^k_{ij}(y)=-c_m\,y_n^{-2}K^k_{ij}$ with the \emph{constant}
  coefficient $K^k_{ij}=-c_i\delta^k_j-c_j\delta^k_i+D^kB_{ij}$
  ($c_i=(\text{finBasis}_i)_n$, $D^k=\sum_\ell B^{k\ell}c_\ell$), and the chart Gram
  \cref{lem:dc-ch8-3-hyperbolic-chart-gram}, $G_{m\ell}=y_n^{-2}B_{m\ell}$ with $B$
  the constant abstract-basis Gram matrix. Every term carries a factor $y_n^{-4}$,
  which cancels against the right-hand side, and the identity reduces to the purely
  finite contraction $\sum_m R^m_{ijk}B_{m\ell}=-(B_{ik}B_{j\ell}-B_{jk}B_{i\ell})$
  in the constants $c,B,D$. Expanding the Christoffel products and contracting each
  index sum with the ``lowered'' identity $\sum_m K^m_{ab}B_{m\ell}=-c_aB_{b\ell}
  -c_bB_{a\ell}+c_\ell B_{ab}$ (which uses $\sum_m D^mB_{m\ell}=c_\ell$), every
  monomial cancels except $-(B_{ik}B_{j\ell}-B_{jk}B_{i\ell})$, the last step using
  the single nontrivial input $\sum_{ij}B^{ij}c_ic_j=1$
  (\cref{lem:dc-ch8-3-inverse-gram-coord}). Hence
  $\langle R(\partial_i,\partial_j)\partial_k,\partial_\ell\rangle_g
  =-(G_{ik}G_{j\ell}-G_{jk}G_{i\ell})$ on the chart frame. Both the pointwise
  curvature form and $-1$ times the standard curvature form
  $\langle x,z\rangle\langle y,t\rangle-\langle y,z\rangle\langle x,t\rangle$ are
  algebraic curvature forms (\cref{lem:dc-ch8-3-pointwise-curvature-form}) that
  agree on this basis frame, so by the basis-determination step of
  \cref{cor:dc-ch4-3-5} they agree on all of $T_pH^n$:
  $\langle R(x,y)z,t\rangle_g=-(\langle x,z\rangle_g\langle y,t\rangle_g
  -\langle y,z\rangle_g\langle x,t\rangle_g)$. This is constant curvature $-1$ at the
  field level; dividing the diagonal $\langle R(x,y)x,y\rangle_g=-|x\wedge y|^2$ by
  $|x\wedge y|^2$ gives $K(p,\sigma)=-1$ for every $2$-plane
  $\sigma=\mathrm{span}\{x,y\}$.
\end{proof}

\begin{lemma}[Christoffel contraction of the hyperbolic metric, closed form]
  \label{lem:dc-ch8-3-christoffel-contraction}
  \lean{Riemannian.Hyperbolic.hyperbolic_chartChristoffelContraction_eq}
  \uses{lem:dc-ch8-3-hyperbolic-christoffel, lem:dc-ch8-3-inverse-gram-coord}
  \leanok
  Contracting the closed-form chart Christoffel symbols
  $\Gamma^k_{ij}=-\delta_{jk}f_i-\delta_{ki}f_j+\delta_{ij}f_k$
  ($f_i=(\partial_i)_n/x_n$) of \cref{lem:dc-ch8-3-hyperbolic-christoffel} against
  two tangent vectors $v,w\in T_pH^n$ collapses, via the inverse-Gram/coordinate
  identity \cref{lem:dc-ch8-3-inverse-gram-coord}, to the basis-independent
  Euclidean-valued expression

  $$
  \Gamma(v,w)(y)=\nabla_v w
    =\frac{1}{y_n}\Big(\langle v,w\rangle\, e_n-v_n\,w-w_n\,v\Big),
  $$

  where $y_n$ is the distinguished coordinate of the base point, $v_n=\langle v,e_n\rangle$,
  $e_n$ is the distinguished unit vector, and $\langle\,,\rangle$ is the ambient
  Euclidean inner product.
\end{lemma}
\begin{proof}
  \leanok
\sketch
  In the abstract chart frame the contraction is
  $\Gamma(v,w)(y)=\sum_k\big(\sum_{ij}\Gamma^k_{ij}(y)\,v^i w^j\big)\partial_k$ with
  $v^i,w^j$ the frame coordinates. Substituting the closed form
  \cref{lem:dc-ch8-3-hyperbolic-christoffel} and summing the two Kronecker terms
  gives $-\big(\sum_i (\partial_i)_n v^i\big)w-\big(\sum_j(\partial_j)_n w^j\big)v
  =-v_n w-w_n v$ (using $\sum_i(\partial_i)_n v^i=v_n$), while the $\delta_{ij}f_k$
  term contracts to $\langle v,w\rangle\sum_k(\sum_\ell B^{k\ell}c_\ell)\partial_k
  =\langle v,w\rangle\,e_n$ by \cref{lem:dc-ch8-3-inverse-gram-coord}; the common
  factor $1/y_n$ is carried throughout.
\end{proof}

\begin{lemma}[geodesic equation of $H^n$ in Euclidean form]
  \label{lem:dc-ch8-3-geodesic-equation}
  \lean{Riemannian.Hyperbolic.hyperbolic_isGeodesic_of}
  \uses{lem:dc-ch8-3-christoffel-contraction}
  \leanok
  A twice-differentiable curve $u:\mathbb R\to\mathbb R^n$ staying in the
  half-space ($u_n(t)>0$) whose acceleration satisfies the ordinary differential
  equation

  $$
  u''(t)+\frac{1}{u_n(t)}\Big(\langle u'(t),u'(t)\rangle\, e_n-2\,u'_n(t)\,u'(t)\Big)=0
  $$

  is a geodesic of $H^n$. This is the geodesic equation $u''+\Gamma(u',u')=0$ read
  through \cref{lem:dc-ch8-3-christoffel-contraction}; because the chart of the open
  half-space is the inclusion, the moving-foot chart curve of the intrinsic geodesic
  predicate is literally $u$.
\end{lemma}
\begin{proof}
  \leanok
\sketch
  For an open subset of $\mathbb R^n$ the extended chart at any base point is the
  inclusion, so the chart-local curve $s\mapsto\varphi_{\gamma t}(\gamma s)$ equals
  $u$ for every $t$; its first and second derivatives are $u'$ and $u''$. The
  intrinsic moving-foot geodesic condition at time $t$ is
  $u''(t)+\Gamma(u'(t),u'(t))(u(t))=0$, and substituting the closed form
  \cref{lem:dc-ch8-3-christoffel-contraction} (with $v=w=u'(t)$, so the two mixed
  terms coincide) gives exactly the displayed equation.
\end{proof}

\begin{lemma}[vertical lines are geodesics of $H^n$]
  \label{lem:dc-ch8-3-vertical-geodesic}
  \lean{Riemannian.Hyperbolic.hyperbolic_vertical_isGeodesic}
  \uses{lem:dc-ch8-3-geodesic-equation}
  \leanok
  For a horizontal base point $c$ (with $c_n=0$) and a constant $a>0$, the affinely
  parametrised vertical line

  $$
  t\longmapsto c+(a\,e^{t})\,e_n,\qquad t\in\mathbb R,
  $$

  perpendicular to the boundary hyperplane $\{x_n=0\}$, is a geodesic of $H^n$. Its
  distinguished coordinate $u_n(t)=a\,e^{t}$ satisfies $u_n''\,u_n=(u_n')^2$.
\end{lemma}
\begin{proof}
  \leanok
\sketch
  Here $u(t)=c+(a e^{t})e_n$, $u'(t)=u''(t)=(a e^{t})e_n$, $u_n(t)=u'_n(t)=a e^{t}$
  and $\langle u',u'\rangle=(a e^{t})^2$. The geodesic equation
  \cref{lem:dc-ch8-3-geodesic-equation} reads
  $(a e^{t})e_n+(a e^{t})^{-1}\big((a e^{t})^2 e_n-2(a e^{t})^2 e_n\big)
  =(a e^{t})e_n-(a e^{t})e_n=0$, which holds identically.
\end{proof}

\begin{lemma}[semicircles perpendicular to $\partial H^n$ are geodesics of $H^n$]
  \label{lem:dc-ch8-3-semicircle-geodesic}
  \lean{Riemannian.Hyperbolic.hyperbolic_semicircle_isGeodesic}
  \uses{lem:dc-ch8-3-geodesic-equation}
  \leanok
  For a center $m$ on the boundary hyperplane ($m_n=0$), a unit horizontal vector
  $\hat u$ ($\hat u_n=0$, $\langle\hat u,\hat u\rangle=1$) and a radius $r>0$, the
  affinely parametrised semicircle

  $$
  t\longmapsto m+(r\tanh t)\,\hat u+(r\,\mathrm{sech}\,t)\,e_n,\qquad t\in\mathbb R,
  $$

  (with $\tanh t=\sinh t/\cosh t$, $\mathrm{sech}\,t=1/\cosh t$), perpendicular to
  the boundary $\{x_n=0\}$ and centered on it, is a geodesic of $H^n$.
\end{lemma}
\begin{proof}
  \leanok
\sketch
  With $u(t)=m+(r\tanh t)\hat u+(r\,\mathrm{sech}\,t)e_n$ one has
  $u'(t)=(r\,\mathrm{sech}^2 t)\hat u-(r\tanh t\,\mathrm{sech}\,t)e_n$; using
  $\mathrm{sech}'=-\tanh\,\mathrm{sech}$, $\tanh'=\mathrm{sech}^2$, and
  $\cosh^2-\sinh^2=1$ one computes $u_n=r\,\mathrm{sech}\,t$,
  $\langle u',u'\rangle=r^2\mathrm{sech}^2 t$, and both the $\hat u$- and
  $e_n$-components of $u''+u_n^{-1}(\langle u',u'\rangle e_n-2u'_n u')$ vanish
  identically. Hence \cref{lem:dc-ch8-3-geodesic-equation} applies.
\end{proof}

\begin{proposition}[geodesics of $H^n$]
  \label{prop:dc-ch8-3-1}
  \lean{Riemannian.Hyperbolic.hyperbolic_geodesic_classification}
  \uses{ex:dc-ch3-3-10, lem:dc-ch8-3-geodesic-equation, lem:dc-ch8-3-vertical-geodesic, lem:dc-ch8-3-semicircle-geodesic}
  \leanok
  \dcref{ch8:3.1}
  The straight lines perpendicular to the hyperplane $x_n=0$, and the circles of
  $H^n$ whose planes are perpendicular to the hyperplane $x_n=0$ and whose centers
  are in this hyperplane, are the geodesics of $H^n$.
\end{proposition}
\begin{proof}
  \leanok
\sketch
  That each stated curve is a geodesic is the forward direction, already established:
  the vertical lines by \cref{lem:dc-ch8-3-vertical-geodesic} and the semicircles
  perpendicular to $\{x_n=0\}$ with centre on it by
  \cref{lem:dc-ch8-3-semicircle-geodesic}, both through the geodesic equation
  \cref{lem:dc-ch8-3-geodesic-equation}. It remains to prove the converse inclusion:
  that these are \emph{all} the geodesics.

  Let $\gamma$ be a geodesic of $H^n$, and read it through the inclusion chart of the
  open half-space, so that $\gamma$ is an ambient curve with $\gamma(0)=p$ and chart
  velocity $\gamma'(0)=v\in\mathbb R^n$. Decompose $v=v_h+v_n\,e_n$, where
  $e_n=\partial/\partial x_n$, $v_n=\langle v,e_n\rangle$, and $v_h$ is the horizontal
  part (so $\langle v_h,e_n\rangle=0$). Through $(p,v)$ passes exactly one member of
  the stated family:
  \begin{itemize}
    \item if $v_h=0$, the reparametrised vertical line
      $s\mapsto c+p_n\,e^{Bs}\,e_n$, with $c=p-p_n e_n$ (so $c_n=0$), $p_n=\langle
      p,e_n\rangle>0$ and $B=v_n/p_n$;
    \item if $v_h\neq0$, the semicircle
      $s\mapsto m+r\tanh(\sigma s+s_0)\,\hat u+r\,\operatorname{sech}(\sigma s+s_0)\,e_n$
      in the plane $\{e_n,\hat u\}$, with $\hat u=v_h/|v_h|$, and centre, radius,
      speed and phase determined by $(p,v)$: writing $s_0=\operatorname{arsinh}(-v_n/|v_h|)$,
      $r=p_n\cosh s_0$, $\sigma=|v_h|\cosh s_0/p_n$ and
      $m=(p-p_n e_n)+\bigl(p_n v_n/|v_h|\bigr)\hat u$, so that $m_n=0$, $\hat u_n=0$,
      $|\hat u|=1$, $r>0$.
  \end{itemize}
  A direct computation (using $\sinh s_0=-v_n/|v_h|$ and
  $\cosh^2-\sinh^2=1$) shows that in either case this family member is a geodesic
  passing through $p$ with chart velocity $v$ at $s=0$. Uniqueness of geodesics with
  prescribed initial data makes $\gamma$ that vertical line or semicircle.
\end{proof}

\section{Space forms}

\begin{theorem}[universal covering of a space of constant curvature]
  \label{thm:dc-ch8-4-1}
  \notready
  \dcref{ch8:4.1}
  \notready
  Let $M^n$ be a complete Riemannian manifold with constant sectional curvature $K$.
  Then the universal covering $\tilde{M}$ of $M$, with the covering metric, is
  isometric to:
  \begin{enumerate}
    \item[(a)] $H^n$, if $K=-1$;
    \item[(b)] $\mathbb{R}^n$, if $K=0$;
    \item[(c)] $S^n$, if $K=1$.
  \end{enumerate}
\end{theorem}
\begin{proof}
  \uses{thm:dc-ch8-2-1,lem:dc-ch8-4-2,prop:dc-ch8-3-const-curv,
    lem:dc-ch7-3-1-covering-diffeo,lem:dc-ch7-3-1-covering-simplyconnected}
\sketch
  $\tilde{M}$ is a simply connected, complete Riemannian
  manifold with constant sectional curvature $K$. Consider first the cases (a) and
  (b) and denote, for convenience, both $H^n$ and $\mathbb{R}^n$ by $\triangle$. Fix
  points $p\in\triangle$, $\tilde{p}\in\tilde{M}$ and a linear isometry
  $i:T_p(\triangle)\to T_{\tilde{p}}(\tilde{M})$. Consider the map
  $f=\exp_{\tilde{p}}\circ\,i\circ\exp_p^{-1}:\triangle\to\tilde{M}$. Since $\triangle$
  and $\tilde{M}$ are complete with non-positive sectional curvature, $f$ is
  well-defined. By \cref{thm:dc-ch8-2-1} it is a local isometry, and
  \cref{lem:dc-ch7-3-1-covering-diffeo} makes it a diffeomorphism.

  For case (c), fix $p\in S^n$, $\tilde{p}\in\tilde{M}$ and a linear isometry
  $i:T_p(S^n)\to T_{\tilde{p}}(\tilde{M})$. Let $q\in S^n$ be the antipodal point of
  $p$ and define $f=\exp_{\tilde{p}}\circ\,i\circ\exp_p^{-1}:S^n-\{q\}\to\tilde{M}$.
  From the Theorem of Cartan, $f$ is a local isometry. Choose $p'\in S^n$,
  $p'\neq p$, $p'\neq q$, set $\tilde{p}'=f(p')$, $i'=df_{p'}$ and define
  $f'=\exp_{\tilde{p}'}\circ\,i'\circ\exp_{p'}^{-1}:S^n-\{q'\}\to\tilde{M}$, where
  $q'$ is the antipodal point of $p'$. On the connected set
  $W=S^n-(\{q\}\cup\{q'\})$ we have $p'\in W$ and $f(p')=\tilde{p}'=f'(p')$,
  $df_{p'}=i'=df'_{p'}$, so by \cref{lem:dc-ch8-4-2}, $f=f'$ on $W$. Define
  $g:S^n\to\tilde{M}$ by $g(r)=f(r)$ if $r\in S^n-\{q\}$ and $g(r)=f'(r)$ if
  $r\in S^n-\{q'\}$. Then $g$ is a local isometry, hence a local diffeomorphism. By
  the compactness of $S^n$, $g$ is a covering map, and since $\tilde{M}$ is simply
  connected, $g$ is a diffeomorphism. Therefore $g$ is an isometry.
\end{proof}

\begin{lemma}[a local isometry pushes geodesics forward]
  \label{lem:dc-ch8-4-2-geodesic-push}
  \lean{Riemannian.Geodesic.isGeodesicOn_comp_of_isLocalDiffeomorph}
  \uses{def:dc-ch1-2-3, def:dc-ch3-2-1}
  Let $f:M\to N$ be a local isometry and let $\gamma$ be a geodesic of $M$ defined on
  an open set $s$ of times. Then $f\circ\gamma$ is a geodesic of $N$ on $s$.
\end{lemma}
\begin{proof}
  \leanok
\sketch
  Being a local isometry, $f$ identifies the metric of $M$ with the pullback $f^*g_N$,
  so $\gamma$ satisfies the $f^*g_N$-geodesic equation. The geodesic equation is a
  pointwise condition, and under $f$ the geodesic operator transforms by the
  differential $df$ alone: from $df\big(\gamma''+\Gamma^{f^*g_N}(\gamma',\gamma')\big)
  =(f\circ\gamma)''+\Gamma^{g_N}\big((f\circ\gamma)',(f\circ\gamma)'\big)$ and the
  vanishing of the left-hand argument one reads off that $f\circ\gamma$ satisfies the
  $g_N$-geodesic equation at each time of $s$.
\end{proof}

\begin{lemma}[the agreement locus is open and closed]
  \label{lem:dc-ch8-4-2-agree-locus}
  \lean{Riemannian.AgreeLocus, Riemannian.mem_agreeLocus_iff, Riemannian.isClosed_agreeLocus, Riemannian.eq_of_pointDatum_of_preconnected}
  Let $f_1,f_2:M\to N$ be differentiable, $M$ connected, and let

  $$
  A=\{x\in M;\ f_1(x)=f_2(x)\ \text{and}\ (df_1)_x=(df_2)_x\}
  $$

  be their \emph{agreement locus}. Then $A$ is closed. If moreover $A$ is open, for
  instance if agreement to first order at a point always propagates to a neighborhood
  of that point, and $A\neq\emptyset$, then $f_1=f_2$.
\end{lemma}
\begin{proof}
  \leanok
\sketch
  A point $x$ lies in $A$ exactly when the tangent maps $df_1$ and $df_2$ agree over
  $x$, since the tangent map of $f$ sends $(x,v)$ to $(f(x),(df)_x v)$: the base
  components give $f_1(x)=f_2(x)$ and the fibre components $(df_1)_x=(df_2)_x$. So $A$
  is the complement of the image, under the bundle projection $TM\to M$, of the locus
  where the two tangent maps differ. The tangent maps are continuous and $TN$ is
  Hausdorff, so that locus is open; the bundle projection is an open map, so its image
  is open and $A$ is closed. If $A$ is also open and nonempty then, $M$ being
  connected, $A=M$, i.e. $f_1=f_2$.
\end{proof}

\begin{lemma}[first-order agreement propagates]
  \label{lem:dc-ch8-4-2-propagation}
  \lean{Riemannian.eventuallyEq_of_pointDatum}
  \uses{lem:dc-ch8-4-2-geodesic-push, prop:dc-ch3-2-9, def:dc-ch3-expmap, lem:dc-ch3-expmap-ray, thm:dc-ch3-2-2}
  \leanok
  Let $f_1,f_2:M\to N$ be local isometries and let $q\in M$ satisfy $f_1(q)=f_2(q)$ and
  $(df_1)_q=(df_2)_q$. Then $f_1=f_2$ on a neighborhood of $q$.
\end{lemma}
\begin{proof}
  \leanok
  \uses{lem:dc-ch8-4-2-geodesic-push,prop:dc-ch3-2-9,def:dc-ch3-expmap,
    lem:dc-ch3-expmap-ray,thm:dc-ch3-2-2}
\sketch
  By \cref{prop:dc-ch3-2-9}, $\exp_q$ carries a ball $B_\varepsilon(0)\subset T_qM$
  diffeomorphically onto an open neighborhood of $q$; so it suffices to prove
  $f_1(\exp_q u)=f_2(\exp_q u)$ for all sufficiently small $u$. Fix such a $u$ and let
  $\gamma_u(t)=\exp_q(tu)$, a geodesic of $M$ on an interval $(-b,b)$ with $b>1$, with
  $\gamma_u(0)=q$ and $\gamma_u'(0)=u$ (\cref{lem:dc-ch3-expmap-ray}). By
  \cref{lem:dc-ch8-4-2-geodesic-push} both $f_1\circ\gamma_u$ and $f_2\circ\gamma_u$ are
  geodesics of $N$ on $(-b,b)$. At $t=0$ they have the same initial point, since
  $f_1(q)=f_2(q)$, and the same initial velocity, since
  $(df_1)_q u=(df_2)_q u$. By uniqueness of geodesics with given initial conditions
  (\cref{thm:dc-ch3-2-2}) they coincide on the connected interval $(-b,b)$; evaluating
  at $t=1$ gives $f_1(\exp_q u)=f_2(\exp_q u)$.
\end{proof}

\begin{lemma}[uniqueness of a local isometry from a point datum]
  \label{lem:dc-ch8-4-2}
  \lean{Riemannian.eq_of_pointDatum_of_localIsometry}
  \uses{def:dc-ch1-2-3}
  \leanok
  \dcref{ch8:4.2}
  Let $f_i:M\to N$, $i=1,2$, be two local isometries of the (connected) Riemannian
  manifold $M$ to the Riemannian manifold $N$, both of dimension $n$ and modelled on a
  common Euclidean space. Suppose that there exists a point $p\in M$ such that
  $f_1(p)=f_2(p)$ and $(df_1)_p=(df_2)_p$. Then $f_1=f_2$.
\end{lemma}
\begin{proof}
  \leanok
  \uses{lem:dc-ch8-4-2-agree-locus,lem:dc-ch8-4-2-propagation}
\sketch
  Consider the agreement locus
  $A=\{x\in M;\ f_1(x)=f_2(x)\ \text{and}\ (df_1)_x=(df_2)_x\}$. It contains $p$, so it
  is nonempty; it is closed by \cref{lem:dc-ch8-4-2-agree-locus}; and it is open,
  because by \cref{lem:dc-ch8-4-2-propagation} each of its points has a neighborhood on
  which $f_1$ and $f_2$ agree, and on such a neighborhood the differentials agree too.
  Since $M$ is connected, $A=M$, that is, $f_1=f_2$.
\end{proof}

The complete constant-curvature manifolds are called \emph{space forms}. An action
of $G$ on $M$ is free and totally discontinuous when $gx=x$ implies $g=e$ and every
$x$ has a neighborhood $U$ with $g(U)\cap U=\varnothing$ for $g\ne e$. Then
$M\to M/G$ is a regular covering.

If $M$ is Riemannian and $\Gamma$ is a totally discontinuous subgroup of its
isometry group, $M/\Gamma$ has a Riemannian structure for which $\pi$ is a local
isometry. For $\tilde p\in\pi^{-1}(p)$ define

$$
\langle u,v\rangle=\langle d\pi^{-1}(u),d\pi^{-1}(v)\rangle_{\tilde{p}},
$$

The definition is independent of the lift because $\Gamma$ acts by isometries.
Completeness and constant curvature descend from $M$ to $M/\Gamma$.

\begin{proposition}[every space form is a quotient of a model]
  \label{prop:dc-ch8-4-3}
  \notready
  \dcref{ch8:4.3}
  \notready
  Let $M$ be a complete Riemannian manifold with constant sectional curvature $K$
  ($1,0,-1$). Then $M$ is isometric to $\tilde{M}/\Gamma$, where $\tilde{M}$ is $S^n$
  (if $K=1$), $\mathbb{R}^n$ (if $K=0$) or $H^n$ (if $K=-1$), $\Gamma$ is a subgroup
  of the group of isometries of $\tilde{M}$ which acts in a totally discontinuous
  manner on $\tilde{M}$, and the metric on $\tilde{M}/\Gamma$ is induced from the
  covering $\pi:\tilde{M}\to\tilde{M}/\Gamma$.
\end{proposition}
\begin{proof}
  \uses{thm:dc-ch8-4-1,lem:dc-ch7-3-1-covering-simplyconnected}
\sketch
  Consider the universal covering $p:\tilde{M}\to M$ with the covering
  metric. Let $\Gamma$ be the group of covering transformations of $p$. Then
  $\Gamma$ is a subgroup of the isometries of $\tilde{M}$ and acts in a totally
  discontinuous manner. Introduce on $\tilde{M}/\Gamma$ the metric induced by
  $\pi:\tilde{M}\to\tilde{M}/\Gamma$. Since the covering $p$ is regular, given
  $\tilde{x},\tilde{y}\in\tilde{M}$, $p(\tilde{x})=p(\tilde{y})$ iff
  $\Gamma\tilde{x}=\Gamma\tilde{y}$, equivalently $\pi(\tilde{x})=\pi(\tilde{y})$.
  The equivalence classes given by $p$ and $\pi$ are the same, which induces a
  bijection $\xi:M\to\tilde{M}/\Gamma$ such that $\pi=\xi\circ p$. Since $\pi$ and
  $p$ are local isometries, $\xi$ is a local isometry, and being a bijection, is an
  isometry of $M$ onto $\tilde{M}/\Gamma$.
\end{proof}

\begin{proposition}[even-dimensional space forms of $K=1$]
  \label{prop:dc-ch8-4-4}
  \notready
  \dcref{ch8:4.4}
  \notready
  Let $M^n$ be a complete Riemannian manifold of even dimension $n=2m$, with constant
  sectional curvature $K=1$. Then $M^n$ is isometric to the sphere $S^n$ or the real
  projective space $P^n$ of the same dimension.
\end{proposition}
\begin{proof}
  \uses{prop:dc-ch8-4-3}
\sketch
  The orthogonal group $O(2m+1)$ is the (transitive) group of isometries of
  $S^n$. Let $\Gamma$ be a subgroup that acts in a totally discontinuous manner on
  $S^n$. If $\gamma\in\Gamma$ has determinant $+1$, there exists an eigenvalue of
  $\gamma$ equal to $1$; then $\gamma$ has a fixed point, which implies $\gamma=e$.
  If $\tilde{\gamma}\in\Gamma$ has determinant $-1$, then $\tilde{\gamma}^2$ has
  determinant $1$, therefore $\tilde{\gamma}^2=e$. The eigenvalues of
  $\tilde{\gamma}$ are $1$ or $-1$. If $1$ is an eigenvalue then $\tilde{\gamma}=e$,
  contradicting $\det\tilde{\gamma}=-1$. As a result $\tilde{\gamma}=-e$, hence
  $\Gamma=\{e,-e\}$. Using \cref{prop:dc-ch8-4-3} we obtain that $M^n$ is either $S^n$, if
  $\Gamma=\{e\}$, or $P^n$, if $\Gamma=\{e,-e\}$.
\end{proof}

\begin{proposition}[metric of constant negative curvature on surfaces of genus $>1$]
  \label{prop:dc-ch8-4-5}
  \notready
  \dcref{ch8:4.5}
  \notready
  Every compact orientable surface of genus $p>1$ can be provided with a metric of
  constant negative curvature.
\end{proposition}
\begin{proof}
\sketch
  In $H^2$ take a closed geodesic polygon $P$ with $4p$ equal sides and identify
  paired sides to form a surface $M^2$. Let $\Gamma$ be
  the subgroup of isometries of $H^2$ generated by the isometries that identify the
  sides of $P$. The translates of $P$ tile $H^2$ iff
  the sum $\alpha$ of the interior angles at the vertices of $P$ equals $2\pi$; in
  this case $M^2=H^2/\Gamma$ has a metric of constant curvature $-1$. We assert that
  it is possible to find a polygon $P$ satisfying $\alpha=2\pi$. By the Gauss–Bonnet
  Theorem, $\alpha=-A+2\pi(2p-1)$, where $A$ is the area of $P$ in the hyperbolic
  metric. If $P$ is arbitrarily small, $\alpha$ is arbitrarily close to $2\pi(2p-1)$;
  if we enlarge the area $A$, $\alpha$ becomes arbitrarily small. By continuously
  deforming $P$, there exists a geodesic polygon such that $\alpha=2\pi$.

\end{proof}

\section{Isometries of the hyperbolic space; Theorem of Liouville}

For a differentiable map $f:U\subset\mathbb R^n\to\mathbb R^n$, conformality at
$p$ is equivalent to the existence of $\lambda(p)>0$ such that

$$
\langle df_p(v_1),df_p(v_2)\rangle=\lambda^2(p)\langle v_1,v_2\rangle,\qquad\lambda^2\neq 0.\qquad(3)
$$

Indeed, if $(3)$ holds then $|df_p(v)|^2=\lambda^2|v|^2$, hence

$$
\cos\angle(df_p(v_1),df_p(v_2))=\cos\angle(v_1,v_2),\qquad(4)
$$

so angles are preserved.

\begin{definition}[conformal map with a coefficient of conformality]
  \label{def:dc-ch8-5-conformal}
  \lean{Riemannian.IsConformalWithCoeff, Riemannian.IsConformalWithCoeff.isConformalMap, Riemannian.IsConformalWithCoeff.conformalAt}
  \leanok
  A map $f:U\subset\mathbb{R}^n\to\mathbb{R}^n$ is \emph{conformal at $p$ with
  coefficient of conformality $\lambda(p)>0$} if $f$ is differentiable at $p$ and
  its differential rescales every inner product by $\lambda^2$, i.e. condition
  $(3)$ holds:
  $\langle df_p(v_1),df_p(v_2)\rangle=\lambda^2(p)\langle v_1,v_2\rangle$ for all
  $v_1,v_2$. Equivalently $df_p$ is a conformal linear map, so $f$ is conformal at
  $p$ in the sense of preserving non-oriented angles.
\end{definition}

\begin{lemma}[a conformal map preserves norms up to $\lambda$ and preserves angles: $(3)\Rightarrow(4)$]
  \label{lem:dc-ch8-5-conformal-angle}
  \lean{Riemannian.IsConformalWithCoeff.norm_fderiv_eq, Riemannian.IsConformalWithCoeff.angle_fderiv_eq}
  \uses{def:dc-ch8-5-conformal}
  \leanok
  If $f$ is conformal at $p$ with coefficient $\lambda$, then $|df_p(v)|=\lambda|v|$
  for every $v$, and $df_p$ preserves the non-oriented angle between any two
  vectors: $\angle(df_p(v_1),df_p(v_2))=\angle(v_1,v_2)$.
\end{lemma}
\begin{proof}
  \leanok
\sketch
  From $(3)$ with $v_1=v_2=v$, $|df_p(v)|^2=\langle df_p(v),df_p(v)\rangle
  =\lambda^2\langle v,v\rangle=\lambda^2|v|^2$, so $|df_p(v)|=\lambda|v|$ since
  $\lambda>0$. Hence the cosine of the angle
  $\langle df_p(v_1),df_p(v_2)\rangle/(|df_p(v_1)|\,|df_p(v_2)|)
  =\lambda^2\langle v_1,v_2\rangle/(\lambda^2|v_1|\,|v_2|)
  =\langle v_1,v_2\rangle/(|v_1|\,|v_2|)$
  is unchanged, giving $(4)$.
\end{proof}

\begin{example}[Basic conformal transformations]
  \label{ex:dc-ch8-5-1}
  \lean{Riemannian.isConformalWithCoeff_isometry, Riemannian.isConformalWithCoeff_dilatation, Riemannian.isConformalWithCoeff_inversion, Riemannian.inversion_unit_involutive, Riemannian.inversion_unit_fixed_of_mem_sphere, Riemannian.inversion_unit_eq_lineMap, Riemannian.dist_inversion_unit_center}
  \uses{def:dc-ch8-5-conformal}
  \dcref{ch8:5.1}
  Euclidean isometries are conformal with coefficient $1$, and dilatations
  $f(p)=\lambda p$, $\lambda>0$, are conformal with coefficient $\lambda$. The
  inversion in the unit sphere centered at $p_0$,
  $$
  f(p)=\frac{p-p_0}{|p-p_0|^2}+p_0,\qquad p\in\mathbb{R}^n-\{p_0\},\qquad(5)
  $$
  is a conformal involution with coefficient $1/|p-p_0|^2$.
\end{example}
\begin{proof}
  \sketch
  The isometry and dilatation cases are immediate. For inversion,
  $$
  df_p(v)=\frac{v|p-p_0|^2-2\langle v,p-p_0\rangle(p-p_0)}{|p-p_0|^4},
  $$
  and direct expansion gives
  $|df_p(v)|^2=|v|^2/|p-p_0|^4$.
\end{proof}

\begin{lemma}[conformal maps are closed under composition]
  \label{lem:dc-ch8-5-conformal-comp}
  \lean{Riemannian.IsConformalWithCoeff.comp, Riemannian.isConformalWithCoeff_id}
  \uses{def:dc-ch8-5-conformal}
  \leanok
  If $f$ is conformal at $p$ with coefficient $\lambda$ and $g$ is conformal at
  $f(p)$ with coefficient $\mu$, then $g\circ f$ is conformal at $p$ with
  coefficient $\mu\lambda$. In particular the identity is conformal with
  coefficient $1$, and any composition of isometries, dilatations, and inversions,
  the class of maps in the Theorem of Liouville, is conformal.
\end{lemma}
\begin{proof}
  \leanok
\sketch
  By the chain rule $d(g\circ f)_p=dg_{f(p)}\circ df_p$. For all $v_1,v_2$,
  $\langle d(g\circ f)_p(v_1),d(g\circ f)_p(v_2)\rangle
  =\mu^2\langle df_p(v_1),df_p(v_2)\rangle
  =\mu^2\lambda^2\langle v_1,v_2\rangle=(\mu\lambda)^2\langle v_1,v_2\rangle$,
  and $\mu\lambda>0$.
\end{proof}

\begin{theorem}[Liouville]
  \label{thm:dc-ch8-5-2}
  \uses{lem:dc-ch8-5-differential-core}
  \notready
  \dcref{ch8:5.2}
  \notready
  Let $f:U\to\mathbb{R}^n$, $n\geq 3$, be a conformal transformation of an open set
  $U\subset\mathbb{R}^n$. Then $f$ is the restriction to $U$ of a composition of
  isometries, dilatations or inversions, at most one of each.
\end{theorem}
\begin{proof}
\sketch
  Let $a_1=(1,0,\dots,0),\dots,a_n=(0,\dots,0,1)$ be the canonical basis and
  $(x_1,\dots,x_n)$ the cartesian coordinates. Let $e_1,\dots,e_n$ be parallel
  differentiable vector fields on $U$ with $\langle e_i,e_j\rangle=\delta_{ij}$. If
  $\lambda$ is the coefficient of conformality of $f$,


  $$
  \langle df(e_i),df(e_k)\rangle=\lambda^2\delta_{ik},\qquad i,k=1,\dots,n.\qquad(6)
  $$


  Let $d^2 f$ be the second differential of $f$, a symmetric bilinear map with
  $d^2 f(a_i,a_j)=\partial^2 f/\partial x_i\partial x_j$. Differentiating $(6)$ with
  $i,j,k$ distinct and combining the three cyclic equations yields
  $\langle d^2 f(e_k,e_j),df(e_i)\rangle=0$ for $i,j,k$ distinct. Letting $i$ vary,
  $d^2 f(e_k,e_j)$ belongs to the plane generated by $df(e_j)$ and $df(e_k)$, so
  $d^2 f(e_k,e_j)=\mu\,df(e_k)+\nu\,df(e_j)$ with $\mu=d\lambda(e_j)/\lambda$,
  $\nu=d\lambda(e_k)/\lambda$, that is,


  $$
  d^2 f(e_k,e_j)=\frac{1}{\lambda}(df(e_k)\,d\lambda(e_j)+df(e_j)\,d\lambda(e_k)).\qquad(7)
  $$


  Put $\rho=1/\lambda$. Using $d(\rho f)=d\rho\,f+\rho\,df$ and $(7)$,


  $$
  d^2(\rho f)(e_k,e_j)=d^2\rho(e_k,e_j)f.\qquad(8)
  $$


  We claim $d^2\rho(e_k,e_j)=0$ for $k\neq j$. Computing the third differential
  $d^3(\rho f)$ via $(8)$ and using the symmetry of the first member in $i,j,k$, one
  obtains $d^2\rho(e_k,e_j)df(e_i)=d^2\rho(e_k,e_i)df(e_j)$; since $df(e_i)$ and
  $df(e_j)$ are linearly independent and $i,j,k$ are distinct, $d^2\rho(e_k,e_j)=0$
  for all $j\neq k$. Choosing $e_1,\dots,e_n$ to form a prescribed orthonormal basis
  at any $p$, this holds for every orthonormal basis; since
  $0=d^2\rho\big(\tfrac{e_j+e_k}{\sqrt 2},\tfrac{e_j-e_k}{\sqrt 2}\big)=\tfrac12\{d^2\rho(e_j,e_j)-d^2\rho(e_k,e_k)\}$,
  we get $d^2\rho(e_j,e_j)=d^2\rho(e_k,e_k)$ for $j\neq k$. In summary, for every
  $p\in U$ and any orthonormal basis, $d^2\rho(e_j,e_k)=\sigma\delta_{jk}$; in the
  canonical basis,


  $$
  \frac{\partial^2\rho}{\partial x_i\partial x_j}=\sigma\delta_{ij}.\qquad(9)
  $$


  Differentiating $(9)$ gives $\partial\sigma/\partial x_i=0$ ($i\neq j$), so
  $\sigma=\text{const.}$

  \emph{Case $\sigma\neq 0$.} Equation $(9)$ implies


  $$
  \rho=\frac{\sigma}{2}\sum x_i^2+\sigma\sum b_i x_i+c,\quad b_i,c\text{ constants},\qquad(10)
  $$


  which we write as $1/\lambda=\rho=a_1|p-p_0|^2+k_1$ with $a_1=\sigma/2$,
  $k_1=\text{const.}$, $p_0\in\mathbb{R}^n$. The proof in this case will be complete
  once we show $k_1=0$: if $k_1=0$, taking the inversion
  $g(p)=\frac{p-p_0}{|p-p_0|^2}+p_0$ and $h=g\circ f^{-1}$, $h$ is conformal with
  coefficient $a_1$, hence an isometry followed by a dilatation; thus
  $f=h^{-1}\circ g$ is an inversion followed by a dilatation, followed by an
  isometry. To show $k_1=0$, apply the argument to $f^{-1}$:
  $\lambda=a_2|f(p)-q_0|^2+k_2$, hence


  $$
  (a_1|p-p_0|^2+k_1)(a_2|f(p)-q_0|^2+k_2)=1.\qquad(11)
  $$


  Equation $(11)$ shows that a sphere of center $p_0$ is mapped by $f$ into a sphere
  of center $q_0$, and since $f$ preserves angles, radii map to radii. For a radial
  segment $p(s)$, $0\leq s\leq s_0$, of arc length $s$,


  $$
  \int_0^{s_0}\Big|df\Big(\frac{dp}{ds}\Big)\Big|\,ds=\int_0^{s_0}\frac{ds}{a_1|p(s)-p_0|^2+k_1}=|f(p(s_0))-f(p(0))|.
  $$


  If $k_1\neq 0$, the left side is a transcendental function of $|p(s_0)-p_0|$, while
  $(11)$ implies it is algebraic. This contradiction shows $k_1=0$, completing the
  case $\sigma\neq0$.

  \emph{Case $\sigma=0$.} Here $\rho=1/\lambda=\sum a_i x_i+c_1=A_1(x)+c_1$. Applying the
  initial argument to $f^{-1}$,


  $$
  (A_1(x)+c_1)(A_2(f(x))+c_2)=1.\qquad(11')
  $$


  Equation $(11')$ shows a hyperplane parallel to $A_1=0$ maps into a hyperplane
  parallel to $A_2=0$; a line perpendicular to $A_1=0$ maps into a line. For a
  segment $p(s)$ of such a line,


  $$
  |f(p(s_0))-f(p(0))|=\int_0^{s_0}\frac{ds}{A_1(p(s))+c_1},
  $$


  which contradicts $(11')$ unless $A_1(p(s))\equiv 0$. We conclude that if
  $\sigma=0$, $\lambda=\text{const.}$; then lengths of tangent vectors are multiplied
  by a constant $\lambda$ and $f$ is an isometry followed by a dilatation.
\end{proof}

\begin{theorem}[isometries of $H^n$ are conformal transformations preserving $H^n$]
  \label{thm:dc-ch8-5-3}
  \uses{def:dc-ch8-5-conformal, def:dc-ch8-3-hyperbolic-space}
  \notready
  \dcref{ch8:5.3}
  \notready
  The isometries of $H^n$ are the restrictions to $H^n\subset\mathbb{R}^n$ of the
  conformal transformations of $\mathbb{R}^n$ that take $H^n$ onto itself.
\end{theorem}
\begin{proof}
  \uses{cor:dc-ch8-2-3,thm:dc-ch8-5-2}
\sketch
  Suppose first $n\geq 3$. Let $f:H^n\to H^n$ be an isometry in the metric
  $g_{ij}$. By Liouville's Theorem, $f$ extends to a conformal map of $\mathbb{R}^n$
  with the metric $\delta_{ij}$. Since $H^n$ is complete in $g_{ij}$, $f$ maps $H^n$
  onto $H^n$. Conversely, let $f:H^n\to H^n$ be a conformal transformation of $H^n$
  onto $H^n$ and let $e_1,\dots,e_n$ be an orthonormal basis (in $g_{ij}$) at
  $p\in H^n$. Since $g_{ij}=\delta_{ij}/x_n^2$ and $f$ is conformal, there exists
  $\lambda^2>0$ with $\langle df_p(e_i),df_p(e_j)\rangle=\lambda^2\delta_{ij}$, so
  $\{df_p(e_i)/\lambda\}$ is orthonormal at $f(p)$. By
  \cref{cor:dc-ch8-2-3} and the global exponential coordinates of $H^n$, there is an isometry $g$ taking
  $p$ to $f(p)$ with $dg(e_i)=df(e_i)/\lambda$. By the first part, $g$ is conformal. Hence
  $h=g^{-1}\circ f$ is the restriction to $H^n$ of a conformal map of $\mathbb{R}^n$
  taking $H^n$ onto $H^n$, leaving $p$ fixed and satisfying $dh_p=\lambda I$. It
  remains to show $h=\text{identity}$.

  Let $P$ be a hyperplane through $p$. By Liouville's theorem, $h(P)$ is a hyperplane
  or a sphere through $p$; since $dh_p$ is a multiple of the identity, $P$ and $h(P)$
  are tangent at $p$. Since $h$ takes $\partial H^n$ into itself and is conformal,
  the angle of $P$ with $\partial H^n$ equals the angle of $h(P)$ with $\partial H^n$.
  We claim $h(P)=P$. Consider a line $r_1$ through $p$ perpendicular to
  $\partial H^n$, $q_1=r_1\cap\partial H^n$. Since $h(r_1)$ is a circle or a line
  making the same angle with $\partial H^n$, $h(r_1)=r_1$, hence $h(q_1)=q_1$. So if
  $P$ is perpendicular to $\partial H^n$, then $h(P)=P$. If $P$ is not perpendicular,
  let $r$ be a line in $P$; for $h(r)$ to be a circle making the same angle
  $\alpha$ it would be necessary that $q_1\in h(r)$, contradicting $h(q_1)=q_1$;
  hence $h(r)$ is a line, $h(r)=r$, and $h(P)=P$. It follows that $h$ cannot be an
  inversion (the image of some plane would be a sphere) or an isometry distinct from
  the identity. From Liouville's theorem, $h$ is a dilatation; since $h$ takes $H^n$
  onto $H^n$, $h$ is the identity.

  For $n=2$, a simple calculation shows the conformal transformations of the form


  $$
  f(z)=\frac{az+b}{cz+d},\quad z\in H^2\subset\mathbb{C},\ a,b,c,d\in\mathbb{R},\ ad-bc>0\qquad(12)
  $$


  (which map $H^2$ onto $H^2$) are isometries of $H^2$ with the metric $g_{ij}$.
  Moreover, for a fixed point $p_0\in H^2$ and unit vector $v_0$ at $p_0$, there
  exists a transformation $(12)$ taking an arbitrary $(p,v)$ to $(p_0,v_0)$ ($p$ and
  $v$ are determined by three parameters, the number of parameters of $f$). Since
  there is a unique orientation-preserving isometry of $H^2$ taking $(p,v)$ to
  $(p_0,v_0)$, all orientation-preserving isometries of $H^2$ are of the form (12).
  Composing these with $z\mapsto-\bar z$ gives the orientation-reversing
  anti-holomorphic isometries, so together they give all isometries of $H^2$.

\end{proof}

\begin{remark}[the ball model; horospheres and hyperspheres]
  \label{rem:dc-ch8-5-4}
  \uses{def:dc-ch8-5-conformal, def:dc-ch8-3-hyperbolic-space, thm:dc-ch8-5-3}
  \notready
  \dcref{ch8:5.4}
  \notready
  Let $B^n=\{p\in\mathbb R^n:|p|<2\}$ carry the metric
  $$
  h_{ij}(p)=\frac{\delta_{ij}}{(1-\tfrac14|p|^2)^2}.
  $$
  It is isometric to $H^n$ via
  $$
  f(p)=4\frac{p-p_0}{|p-p_0|^2}-(0,\dots,1),\qquad p_0=(0,\dots,-2),\qquad(13)
  $$
  A map $g:B^n\to B^n$ is an isometry for $h_{ij}$ if and only if it is the
  restriction of a conformal transformation of $\mathbb{R}^n$ taking $B^n$ onto
  $B^n$.
  Euclidean spheres contained in $H^n$ are geodesic spheres. Spheres tangent to
  $\partial H^n$ are intrinsically flat \emph{horospheres}. Spheres meeting
  $\partial H^n$ are \emph{equidistant hypersurfaces}: they lie at fixed normal
  distance from a totally geodesic hyperplane.
\end{remark}
\begin{proof}
  \sketch
  The identities
  $|df_p(v)|^2=16|v|^2/|p-p_0|^4$ and
  $f_n(p)=(4-|p|^2)/|p-p_0|^2$ show that (13) pulls back the half-space metric to
  $h$. Conjugating by (13) gives the isometry classification. Inversions send
  the three types of Euclidean spheres to centered spheres, boundary-parallel
  hyperplanes, and hyperplanes at fixed distance from a totally geodesic one,
  respectively.
\end{proof}

\begin{lemma}[a uniform radial chart neighborhood]
  \label{lem:dc-ch8-2-1-radial-chart}
  \dcref{ch8:2.1}
  \lean{Riemannian.Jacobi.exists_radial_chart_neighborhood,
    Riemannian.Jacobi.exists_pair_radial_chart_neighborhood}
  \leanok
  Let $p$ be a point of a complete Riemannian manifold and let $\exp_p$ be defined on the
  whole tangent space. There is an open neighborhood $W$ of $0\in T_pM$ such that every
  radial geodesic $t\mapsto\exp_p(tw)$ with $w\in W$ stays in the source of one fixed chart
  at $p$ for $0\leq t\leq1$. For two complete manifolds modelled on the same space $E$ and a
  continuous linear equivalence $i:E\to E$ between their tangent coordinates, one may choose
  the same $W$ for the two radial families, using the intersection with the inverse image of the
  second neighborhood under $i$.
\end{lemma}
\begin{proof}
  \leanok
  The global exponential map is continuous and sends $0$ to $p$, while the chart source is
  open and contains $p$. Its inverse image therefore contains a metric ball about $0$; radial
  homogeneity and $0\leq t\leq1$ keep every scaled vector in that ball. For a pair, intersect
  the first ball with the pullback of the second ball under $i$.
\end{proof}
